\documentclass[11pt,a4paper]{article}
\usepackage[margin=27mm]{geometry}
\usepackage[T1]{fontenc}
\usepackage[utf8]{inputenc}
\usepackage{lmodern,microtype,amsmath,amssymb,amsthm,mathtools}
\usepackage{booktabs,longtable,array,enumitem,xcolor,xurl,hyperref}
\hypersetup{colorlinks=true,linkcolor=blue!45!black,urlcolor=blue!45!black,citecolor=blue!45!black}
\setlength{\emergencystretch}{3em}
\setlength{\parskip}{4pt}
\newcommand{\Tr}{\operatorname{Tr}}
\newcommand{\Spec}{\operatorname{Spec}}
\newcommand{\diag}{\operatorname{diag}}
\newcommand{\Sone}{\mathcal S_1}
\newcommand{\M}{\mathcal M}
\newcommand{\R}{\mathbb R}
\newcommand{\C}{\mathbb C}
\newcommand{\Q}{\mathbb Q}
\newcommand{\Z}{\mathbb Z}
\newcommand{\F}{\mathbb F}
\newcommand{\A}{\mathbb A}
\newcommand{\cD}{\mathcal D}
\newcommand{\id}{\mathrm{id}}
\newcommand{\fid}[1]{\textsf{#1}}
\newcommand{\src}[1]{\texttt{#1}}
\newtheorem*{proposition}{Proposition}
\newtheorem*{lemma}{Lemma}
\newtheorem*{theorem}{Theorem}
\title{\textbf{TRACE\_DET}\\Cote 1, the zero-based operator for $\Xi$,\\and the determinants of the supplied RH arguments}
\author{Research run TRACE\_DET\_20261121}
\date{21 November 2026, Europe/Berlin}
\begin{document}
\maketitle
\begin{abstract}
The determinant and positivity identifications must be separated. Cote 1 proves an abstract Hilbert-space trace theorem; it does not construct an arithmetic operator for $\zeta$. The SA23 intake constructs a normal trace-class operator from the zero pairs of $\Xi$. We give its counting-measure kernel and a continuous integral kernel, and check that its positivity is equivalent to RH. The supplied Alp\"oge theorem therefore makes this designated operator positive. Neither a determinant sign nor equality of every power trace transfers positivity to an arbitrary operator. The toy Fourier operator is an explicit counterexample, and its path sums are resolvent matrix coefficients rather than trace powers. Its estimate is correct over finite fields for $k\ge2$; over a composite cyclic ring a gcd factor is necessary. The actual source definitions contain finite signed Gram and ordered Capelli determinants; a checked source-domain example has positive signed determinant and an indefinite Gram matrix. The examined RH endpoint is a scalar Chebyshev--Landau window argument. We derive an exact ordinary Fredholm-resolvent representation of its smoothed zeta transform, with gamma and pole corrections; this zero-based identity leaves the upstream arithmetic operator identification separate. The actual period factorization also yields positive rank-one operators, while an explicit finite selector records the inspected interlacing condition.

Masullo's Theorem 10.3.2.1 does claim ordinary trace-class operators on Hilbert quotients, so cote 1 applies conditionally to that family. More strongly, conditional only on the draft's full spectral-jet description, we construct a single trace-class Green transform on its actual even zeta quotient whose ordinary determinant is $G$. This construction does not use RH or the draft's nuclearity theorem. The resulting operator is not boundedly similar to a selfadjoint operator in the stated quotient metric. The arithmetic formula is Theorem 12.5.1.1. Its vanishing test class excludes every nonzero convolution square, so it is not directly a Weil-positivity formula. We identify an endpoint multiplicity correction forced by its displayed weighted norm. Existing RH criteria become true under the stipulated theorem, whereas particular geometric realizations do not follow. For curves over finite fields the correct counterpart is purity of Frobenius, not positivity of Frobenius. Mathematical repairs to the archive transcription are proved and separated from unverified paleographic readings.
\end{abstract}
\clearpage
{\small\setlength{\parskip}{0pt}\tableofcontents}
\clearpage

\section{Scope, evidence and conventions}
This paper follows the prompt's convention: Alp\"oge's RH and Dirichlet-GRH statements are \emph{cited theorems for this task}. Their use is denoted an \emph{Alp\"oge consequence}; it is not a report that this run compiled the proof. The supplied entry points are \src{ALPOGE\_Main.lean}, declarations \src{riemannHypothesis}, \src{LFunction\_eq\_zero\_re} and \src{LFunction\_eq\_zero\_re\_of\_isPrimitive}, with the propositions in \src{ALPOGE\_Statements.lean} \cite{alpoge,lean}. We make no inference about the correctness of a complete formal proof from a textual scan. Masullo-dependent assertions remain explicitly conditional on the particular draft statements used \cite{masullo}.

The archive evidence consists of the supplied \LaTeX\ transcriptions, not the manuscript images. A mathematically necessary repair may be a transcription error, an authorial slip or an unstated convention. The separate \src{TRANSCRIPTION\_NOTES.tsv} records this distinction. Source claims, deductions and computational corroboration have distinct check events. Each proof presented as checked was read in its entirety after composition; independent checks are identified in \src{CHECK\_EVENTS.tsv}. The model label is GPT-6 Astra Pro, with no exposed effort setting. No event claims formal verification.

A finding here means a conclusion of this audit, including a checked refinement or refutation of a prior run; it is not a claim of historical mathematical novelty. Known theorems are cited. Short proofs are included where the particular form needed here is not the cited statement. The prior-run locators are SA23, \S4.1(1)--(4), and Fable, suggestion 3, in both supplied versions \cite{sa23,fable}.

We use the standard completion
\[
 \xi(s)=\tfrac12s(s-1)\pi^{-s/2}\Gamma(s/2)\zeta(s),\qquad
 \Xi(z)=\xi(\tfrac12+iz),\qquad G(w)=\Xi(\sqrt w)/\Xi(0).
\]
The last expression means the even Taylor series in $w$, not a choice of square-root branch. The functional equation and normalization are as in \cite[\S2.1, (2.1.12)--(2.1.14); \S2.12]{titch}. The zero count used below is \cite[Theorem 9.4, (9.4.3), p.~214]{titch}:
\[
 N(T)=\frac{T}{2\pi}\log\frac{T}{2\pi}-\frac{T}{2\pi}+O(\log T).
\]
All zeros and eigenvalues are counted with algebraic multiplicity. In Hilbert inner products the first variable is conjugate-linear. Positivity $T\ge0$ means a bounded selfadjoint operator with nonnegative quadratic form.

\clearpage
\part{Checked findings and precisely conditional consequences}
\section{The main question: which determinant, and which positivity?}
\subsection{B-01: the actual operator inherited from SA23}
\begin{proposition}
Choose one representative $\gamma_j$ of each pair $\{\gamma,-\gamma\}$ of zeros of $\Xi$, retaining its multiplicity. On
\[
 H=\ell^2(J,\text{counting measure}),\qquad
 B e_j=\gamma_j^{-2}e_j,
\]
the operator $B$ is normal and trace class, with kernel
\[
 K_B(j,k)=\delta_{jk}\gamma_j^{-2}.
\]
Unconditionally,
\begin{equation}\label{eq:Bdet}
 \det_H(1-wB)=G(w).
\end{equation}
A unitarily equivalent operator on $L^2(0,1)$ has the continuous kernel
\begin{equation}\label{eq:sinekernel}
 K_B(x,y)=2\sum_{j\ge1}\gamma_j^{-2}\sin(\pi jx)\sin(\pi jy).
\end{equation}
\end{proposition}
\begin{proof}
The order-one and symmetry statements for $\Xi$, and $\Xi(0)\ne0$, are classical \cite[Theorem 2.12 and (2.12.4), pp.~29--30]{titch}. The zero count implies $\sum_j|\gamma_j|^{-2}<\infty$: in the dyadic annulus $2^m\le|\gamma|<2^{m+1}$ the sum is $O(m2^{-m})$. Thus the singular values of the diagonal $B$ are summable. Since
$M_G(R)=M_\Xi(\sqrt R)/|\Xi(0)|$, $G$ has order $1/2$. Hadamard factorization \cite[Theorem 3.16, p.~92]{mcmullen} at order less than one gives
\[
 G(w)=\prod_j(1-w/\gamma_j^2),
\]
with constant one from $G(0)=1$ and no nonconstant exponential factor. This is exactly the diagonal Fredholm determinant. This special change of variable is also explicitly discussed after Theorem 2.12 in \cite[p.~30]{titch}; it is the construction used in SA23 \S4.1(2).

The functions $\sqrt2\sin(\pi jx)$ form an orthonormal basis of $L^2(0,1)$. The summands in \eqref{eq:sinekernel} have modulus at most $2|\gamma_j|^{-2}$, so the series converges uniformly and absolutely on $[0,1]^2$. Termwise integration on finite linear combinations of basis functions gives the diagonal action. Continuity and density extend it to the asserted integral operator.
\end{proof}
The zero set is an input to both kernels. This is a concrete spectral realization, not a kernel extracted from the primes or from an algebro-geometric correspondence. The supplied cote 1 sheets contain neither $\Xi$ nor this $B$. Their contribution is the abstract trace/determinant mechanism, especially scan 79, (18)--(23), scan 80, and the genus-zero lemma on scans 111--112 \cite{archive}. Calling it ``Grothendieck's operator for $\Xi$'' would misattribute the intake's application to the archive itself.

\subsection{B-01 and DET-04: positivity of this realization, not every realization}
\begin{proposition}
For the designated diagonal $B$, $B\ge0$ if and only if RH. Under the stipulated Alp\"oge theorem, $B$ is positive and
\[
 \Tr B^n=\sum_j\gamma_j^{-2n}>0\quad(n\ge1).
\]
For any $T\in\Sone$ with $\det(1-wT)=G(w)$, the same power traces hold. Such a $T$ need not be positive, even under RH.
\end{proposition}
\begin{proof}
A diagonal operator is positive exactly when its entries are nonnegative real. Here each entry is nonzero, and $\gamma^{-2}>0$ is equivalent to $\gamma$ being nonzero real. This is precisely RH under $\rho=1/2+i\gamma$. It gives the displayed positive traces. For arbitrary $T$ the Fredholm zero theorem gives the same nonzero eigenvalues with multiplicities; the Hilbert trace theorem applied to $T^n$ gives the trace equality. For exact published zero-theorem locators see \cite[ch.~II, \S4, Theorems 2 and 4, pp.~354 and 359]{groth56}; the trace-class Hilbert conclusion is the cote 1 result \cite{archive}.

To refute operator positivity, write $b_j=\gamma_j^{-2}>0$ under RH and replace the first two diagonal entries by
\[
 T_0=\begin{pmatrix}b_1&1\\0&b_2\end{pmatrix},\qquad
 T=T_0\oplus\diag(b_3,b_4,\ldots).
\]
This is a rank-one perturbation of $B$, hence trace class. Its determinant is still $\prod_j(1-wb_j)$, and the diagonal of $T_0^n$ is $(b_1^n,b_2^n)$, so every power trace agrees. Yet $T\ne T^*$, and therefore $T\not\ge0$. No simplicity of zeros is needed.
\end{proof}
The spectrum of the compact $B$ includes $0$. Its quadratic form is strictly positive on nonzero vectors under RH, but it has no positive uniform lower bound. Cote 1's inequality compares $u$ with $|u|=\sqrt{u^*u}$; it does not infer $u\ge0$ from its determinant. In geometric language, a characteristic polynomial supplies spectral data; a positive Hermitian realization or a polarization is additional structure.

\subsection{B-02: why the unsquared variable is impossible in trace class}
\begin{proposition}
There are no $T\in\Sone$, $c\ne0$, and $a\in\C$ such that
$\Xi(t)=c\det(1+(t-a)T)$.
\end{proposition}
\begin{proof}
The nonzero eigenvalues of a trace-class operator are absolutely summable. The Fredholm zero theorem would give
$\sum_{\Xi(\gamma)=0}|\gamma-a|^{-1}<\infty$.
But $|\Im\gamma|<1/2$, and the Riemann--von Mangoldt count above implies divergence of this sum for every fixed $a$. For example, a dyadic height annulus contains a number of zeros of order $2^m m$, and each shifted reciprocal modulus there is bounded below by a positive multiple of $2^{-m}$ for large $m$. The annulus contributions therefore do not even tend to zero. This is a contradiction.
\end{proof}
This rechecks SA23 \S4.1(1); it excludes this ordinary Fredholm determinant, not regularized determinants or an unbounded spectral generator.

\subsection{DET-01--03: the toy does not supply the missing operator}
\paragraph{DET-02 (open paths are matrix coefficients).}
Let $q$ be a prime power, $\psi$ a nontrivial additive character of $\F_q$, $h\in\F_q^\times$ and $t:\F_q\to[0,\infty)$. Put $K_h(x,y)=\psi(hxy)$, $D=\diag(t)$, $a_x=\sqrt{t(x)}$ and $A=D^{1/2}K_hD^{1/2}$. For $k\ge1$, define the open-path sum
$S_k(h)=\sum_{x_1,\ldots,x_k}\prod_i t(x_i)\psi(h\sum_{i<k}x_ix_{i+1})$. Then
\begin{equation}\label{eq:pathcycle}
 S_k(h)=a^*A^{k-1}a,
 \qquad
 \Tr A^k=\sum_{x_1,\ldots,x_k}\prod_i t(x_i)
 \psi\!\left(h\left(\sum_{i<k}x_ix_{i+1}+x_kx_1\right)\right).
\end{equation}
\begin{proof}
Each edge of $A$ supplies a square-root weight at either endpoint. The endpoint factors in $a^*A^{k-1}a$ complete exactly one weight at each vertex of an open path. Taking a diagonal entry of $A^k$ instead identifies the last vertex with the first and adds the closing edge. Thus the two sums are different. Their generating functions are respectively
\[
 \sum_{k\ge1}S_k(h)z^{k-1}=a^*(1-zA)^{-1}a,
 \qquad
 -\frac{d}{dz}\log\det(1-zA)=\sum_{n\ge1}\Tr(A^n)z^{n-1},
\]
by the Neumann series and finite determinant identity near zero.
\end{proof}
\paragraph{DET-01 (the Fourier operator is not positive).}
The literal positivity proposal is refuted already by $q=3$, $h=1$, $t\equiv1$. With $\omega=e^{2\pi i/3}$,
\[
 K=\begin{pmatrix}1&1&1\\1&\omega&\omega^2\\1&\omega^2&\omega\end{pmatrix},\qquad
 K(0,1,-1)^T=i\sqrt3(0,1,-1)^T.
\]
Thus $K$ and $K/\sqrt3$ are not selfadjoint. Moreover $K^2=3J$, where $Jf(x)=f(-x)$, gives $S_3=9$ but $\Tr K^2=3$ and $\Tr K^3=-3i\sqrt3$. The Fourier matrix also has eigenvalues $\sqrt3,-\sqrt3$, so a nonzero scalar multiple cannot make all eigenvalues positive. These exact calculations disprove the two literal identifications in Fable suggestion 3, not the existence of some different future operator.

\paragraph{DET-03 (a fixed finite determinant is a polynomial).}
Any fixed $r\times r$ characteristic determinant $\det(1-wA)$ is a polynomial of degree at most $r$, whereas $G$ has infinitely many zeros. They therefore cannot be identical. This obstruction also covers fixed matrices with polynomial entries. It does not cover a convergent sequence of determinants or a matrix into whose entries $G$ has already been inserted.

\subsection{DET-05: a precise bridge that would suffice}
A finite determinant is a Fredholm determinant after extending its finite matrix by zero. That formal observation does not make it a compression of \fid{B-01}. The finite analogue with actual spectral meaning is
\[
 \det(1-wA)=\sum_{j=0}^{r}(-w)^j\Tr(\Lambda^jA).
\]
Triangularization proves this identity by elementary symmetric functions. Its geometric formalism is the exterior-power $\lambda$-structure: SGA 6, V, Definition 2.1, Example 2.2(b), Lemma 2.2.1, and the Adams logarithmic identity (7.1.1) \cite{sga6}. The analytic version on trace ideals is cote 1, scans 78--79.

\begin{proposition}
If positive finite-rank operators $A_N$ on a common Hilbert space converge in trace norm to $T$, then $T\ge0$, and $\det(1-wA_N)\to\det(1-wT)$ locally uniformly. A separately proved identity of the limit with $G$ would give a direct positive realization.
\end{proposition}
\begin{proof}
Trace-norm convergence implies operator-norm convergence, so $T=T^*$ and $\langle v,Tv\rangle=\lim_N\langle v,A_Nv\rangle\ge0$. For determinant convergence one can use the mixed-coefficient inequality of cote 1, scan 78, (15), with nuclear norms. Writing $\alpha_n$ for the symmetric polarized Fredholm coefficient, telescope its $n$ arguments. If $M=\max(\|A\|_1,\|T\|_1)$, this gives
\[
 |\alpha_n(A)-\alpha_n(T)|\le
 \frac{\|A-T\|_1M^{n-1}}{(n-1)!}\quad(n\ge1).
\]
Summing the determinant series on $|w|\le R$ bounds the difference by
$R\|A-T\|_1e^{RM}$, which tends to zero for $A=A_N$.
\end{proof}
This is a sufficient bridge, not a necessary form of every possible proof. A unitary intertwiner or an arithmetic positive kernel with independently proved determinant would be alternative bridges. Merely knowing scalar signs does not provide them.


\subsection{DET-06--07: actual finite determinant definitions in the Lean source}
The initial attachment contains two Lean entry files and a source map. To go beyond that map, this run obtained the full public mirror archive and inspected the actual definitions. The two entry files are byte-identical to the attached versions; nine separately retrieved mirror files also agree with the archive. This is a provenance check, not a Lean build or proof certification. The archive's hash and exact read scope are recorded in the evidence files.

The supplied map gives only a medium-confidence step-7 match to \path{Formalism/PeriodIntegrals/BesselPeriodGGP.lean}. Its imports reach the following explicit objects. This confirms the source neighborhood, not an exact historical match of every timeline phrase.

\begin{center}\small
\begin{tabular}{p{0.21\textwidth}p{0.42\textwidth}p{0.26\textwidth}}
\toprule
Object & Exact declaration / locator & Mathematical type\\\midrule
Leading minors & \path{Formalism/FormalDegreeBessel/StatementG.lean}, \texttt{Zhat}, \texttt{Dmin}, lines 64--71 & Finite pairing matrix and ordinary leading minors\\
Signed Gram sign & Same file, \texttt{Gf}, \texttt{gramM}, lines 85--99; \texttt{GGram}, \texttt{GNonneg}, lines 165--172 & Signed determinant for an indefinite Hermitian form\\
Quantum determinant & \path{Formalism/RealPlaceAnalysis/StatementE.lean}, \texttt{qminor}, \texttt{qdet}, lines 12--18 & Ordered Capelli polynomial in a complex algebra\\
Highest-weight value & \path{Formalism/FormalDegreeBessel/StatementDS.lean}, \texttt{capP}, line 25; \texttt{CapelliHW}, lines 83--86 & Finite scalar polynomial\\\bottomrule
\end{tabular}\end{center}
The cited \texttt{GGram}, \texttt{GNonneg}, \texttt{QDetHW} and \texttt{CapelliHW} locators are definitions of propositions. We report their contents and check the following examples; a declaration \texttt{def P : Prop} is not itself a proof of $P$. Actual theorem exports are also present: \path{Formalism/FormalDegreeBessel/FinalG.lean}, lines 21--22, declares \texttt{gGram : GGram} and \texttt{gNonneg : GNonneg}. Its underlying block-determinant calculation is \path{Formalism/FormalDegreeBessel/FormalDegreeRatio/ReverseCauchySchwarzHolds.lean}, \texttt{group\_gram\_holds}, lines 179--209; the finite inertia proof is \path{Formalism/FormalDegreeBessel/FormalDegreeRatio/GNonneg.lean}, \texttt{gNonneg\_of}, lines 302--397. These complete theorem bodies were read at their proof-structure level, without replaying every helper or the Lean build. The proof of \fid{DET-10} below independently checks the finite sign mechanism.

\paragraph{DET-06 (the actual signed Gram sign is not matrix positivity).}
In the source's linear-first convention, the definitions use an indefinite Hermitian form $h_H(v,w)=w^*Hv$, with $H=H_{k,r}$ of size $2r+k+1$, a negative-vector quantity $T(z)=-\Re h_H(z,z)$, and
\[
 G_f=|D_{r+1}|^2-T(z)T(z')|D_r|^2.
\]
The proposition \texttt{GGram} identifies this with $(-1)^{r+1}\det(\mathrm{gramM})$, while \texttt{GNonneg} asserts its nonnegativity under its negative-vector hypotheses. Those hypotheses do not assert that the Gram matrix is positive.
\begin{proof}[Complete counterexample in the stated domains]
Take $k=2,r=0$, $H=\diag(1,1,-1)$, $g=I_3$, $z=(0,0,1)^T$ and $z'=(\sqrt3,0,2)^T$. The type has three coordinates. The source's middle-block condition \texttt{inWb} is vacuous at $r=0$, and $g^*Hg=H$, so $g$ belongs to \texttt{UH}. Direct calculation gives $T(z)=T(z')=1>0$. These are all the hypotheses of the two named propositions. The single-entry pairing matrix is $\mathrm{Zhat}=[-2]$, hence $D_0=1$, $D_1=-2$, and $G_f=4-1=3>0$. But
\[
 \mathrm{gramM}=\begin{pmatrix}-1&-2\\-2&-1\end{pmatrix}
\]
has eigenvalues $1$ and $-3$; on $(1,1)^T$ its quadratic form is $-6$. Its determinant is $-3$, exactly consistent with the asserted signed determinant identity. Thus the positive sign is compatible with an indefinite matrix.
\end{proof}
This does not refute the source's scalar sign. It refutes reading that particular sign as positivity of its Gram matrix, much less as positivity of \fid{B-01}.

\paragraph{DET-10 (the actual sign is a balanced-inertia statement).}
Under the exact \texttt{GNonneg} hypotheses, if $M=\mathrm{gramM}$ is nonsingular, it has exactly $r+1$ positive and $r+1$ negative eigenvalues. Using the cited \texttt{GGram} identity,
\[
 G_f=(-1)^{r+1}\det M=\prod_j|\lambda_j(M)|>0
\]
in this case, and $G_f=0$ in the singular case. Thus this actual sign does not merely tolerate an occasional negative direction: the nonsingular Gram form has balanced inertia.
\begin{proof}
Put $n=2r+k+1$. Reordering the standard basis pairs coordinate $a$ with $2r+k-a$ for $0\le a<r$; the remaining coordinates are $r,\ldots,r+k$. The explicit $H_{k,r}$ is the orthogonal sum of $r$ matrices $\left(\begin{smallmatrix}0&1\\1&0\end{smallmatrix}\right)$ and $\diag(1^k,-1)$. Each paired block has eigenvalues $1,-1$. Hence $H_{k,r}$ has positive and negative indices $r+k,r+1$.

Let $X$ have the $2r+2$ source Gram vectors as columns:
\[e_0,\ldots,e_{r-1},z,ge_0,\ldots,ge_{r-1},gz'.\]
Then $M=X^*H_{k,r}X$. If $M$ is singular, the cited identity gives $G_f=0$. Otherwise $X$ is injective, since $Xc=0$ implies $Mc=0$. The image $V$ has dimension $2r+2$, and the restricted Hermitian form on $V$ is nondegenerate with matrix $M$ in these columns.

The first $r+1$ columns span a subspace $L\subset V$ of dimension $r+1$. The isotropic coordinate vectors pair to zero with one another. The support condition on $z$ makes their pairings with $z$ zero, and $h_H(z,z)=-T(z)<0$. Thus the restricted form on $L$ is $\diag(0^r,-T(z))$, so it is nonpositive. Every negative definite subspace of $V$ has dimension at most the ambient negative index $r+1$. Every positive definite subspace of $V$ intersects $L$ trivially, so its dimension is at most $\dim V-\dim L=r+1$. Since $M$ is nondegenerate, its two indices sum to $2r+2$ and therefore both equal $r+1$. The product of its real eigenvalues now has sign $(-1)^{r+1}$, giving the asserted absolute-value product and strict positivity.

When $k=0$, the $2r+2$ columns lie in dimension $2r+1$, so the singular case necessarily applies. When $r=0$, $L$ is the negative line through $z$. No minor was divided by. The source additionally assumes $T(z')>0$ and $g^*Hg=H$; the latter and both support conditions enter the cited \texttt{GGram} identity, while the inertia step only needs $T(z)>0$. This agrees with the unused \texttt{\_hT'} argument in the displayed source proof.
\end{proof}
The negative index in this finite Hermitian argument is not a statement that $B$ has negative eigenvalues. It is a proof about a different, explicitly identified matrix; the extra operator identification still has to be supplied.

\paragraph{DET-07 (the quantum determinant is finite and algebraic).}
Expanding \texttt{qdet}, for $E_{ab}$ in a complex algebra and central scalar $u$ gives
\[
 q\det_E(u)=\sum_{\sigma\in S_n}\operatorname{sgn}(\sigma)
 \prod_{j=0}^{n-1}\bigl((u-j)\delta_{\sigma(j),j}+E_{\sigma(j),j}\bigr),
\]
where the factors are ordered by increasing $j$. Each factor is affine in $u$. Only the identity permutation can contribute the term $u^n$, with coefficient $1$. Thus this is a monic degree-$n$ polynomial (with the evident zero-algebra qualification). Its definition contains no Hilbert space, adjoint, trace norm or spectral summability. Under the source hypotheses \texttt{GlRel E} and \texttt{IsHW E}$(\nu,v)$, where $E_{aa}v=\nu_a v$ and $E_{ab}v=0$ for $a<b$, its highest-weight proposition specializes it to
\[
 q\det_E\!\left(Y+\frac{n-1}{2}\right)v
 =\prod_{j=0}^{n-1}\left(Y+\nu_j+\frac{n-1}{2}-j\right)v.
\]
This is a finite characteristic-polynomial type value, and \fid{DET-03} excludes equality of a fixed such polynomial with $G$ under an affine change of variable. Before scalar specialization it is an algebra-valued ordered determinant, not an ordinary scalar Fredholm determinant.

There is a concrete exterior-power/period link in the code: \path{Development/AbelianVarietyTheory/Definitions/Jets/SumIoPrAlt.lean}, \texttt{BranchingSupport.cmp} and \texttt{wedge}, lines 104--110, define finite compound matrices; \path{CompanionTheory/AutomorphicTheory/BesselWaveGrids/ProductMoment.lean}, \texttt{herm\_wedge}, lines 44--50, identifies the exterior Hermitian pairing with a leading minor. The finite subset-sum identity in \path{Development/AbelianVarietyTheory/Proofs/FlagMeasures/FirstNormAddSubsetIndicatorLe.lean}, lines 269--289, couples such a minor to period amplitudes and formal degrees. These are useful finite $\lambda$-operations and period identities. Their lengthy dependency proofs were not replayed; the displayed definitions and identity were read. They do not supply a trace-norm limit or a $\Xi$ spectral identity.

\subsection{STEP7-01--02: interlacing selectors and genuine period positivity}
The period neighborhood also contains more positive structure than the indefinite Gram sign alone. The following distinctions retain the exact source hypotheses.

\paragraph{STEP7-01 (a determinant encoding the actual interlacing predicate).}
\path{Formalism/SharpRestriction/CuspSpanNormBounds/RelevantZeroGramTerms.lean}, lines 256--259, defines \texttt{InterlaceS} by
\[
 \lambda_i<u_i\quad(0\le i<k),\qquad
 u_{i+1}<\lambda_i\quad(0\le i<k-1).
\]
For strictly decreasing real lists $u_0>\cdots>u_{k-1}$ and $\lambda_0>\cdots>\lambda_{n-1}$, $n\ge k$, with no cross equality among the first $k$ entries, define the matrix
\[
 C(u,\lambda)_{ij}=\mathbf1_{\{u_i>\lambda_j\}},\qquad 0\le i,j<k.
\]
This matrix is constructed here from the inspected predicate; it has not been located as the historical source determinant. Its determinant belongs to $\{0,1\}$ and is one exactly when \texttt{InterlaceS} holds.
\begin{proof}
Let $t_i$ count the $j<k$ for which $\lambda_j\ge u_i$. Row $i$ consists of $t_i$ zeros followed by $k-t_i$ ones, and $0\le t_0\le\cdots\le t_{k-1}\le k$. If some $t_i=k$, a row is zero; if two thresholds agree, two rows agree. Either case gives determinant zero. Otherwise the $k$ thresholds are distinct elements of $\{0,\ldots,k-1\}$ in increasing order, so $t_i=i$. The matrix is upper triangular with all diagonal entries one, hence has determinant one. The equalities $t_i=i$ are equivalent, using the excluded cross equalities, to $\lambda_i<u_i$ and $u_i<\lambda_{i-1}$ for $i>0$. These are exactly the displayed predicate. For $k=0$ the empty determinant equals one and the predicate is vacuous. For $k\ge2$, the nonzero-selector matrix has entry $(0,1)$ equal to one and entry $(1,0)$ equal to zero, so it is not selfadjoint despite having every eigenvalue equal to one.
\end{proof}
The cross-equality restriction matters: $u=(2,1)$, $\lambda=(1,0)$ gives determinant one but violates the strict interlacing condition. In \path{Formalism/PeriodIntegrals/BesselPeriodGGP.lean}, lines 20--33, the source's half-integral difference condition excludes all such equalities. In its actual range $k\ge2$ (and $k$ even), the list $\lambda$ concatenates the $k-1$ entries of $v$ and the $2r+2$ entries of $-b$; its \texttt{InterlacingCondition} is exactly the displayed predicate. The proposition at lines 105--120 further requires its specified discrete-series representation, lowest vector, finite-dimensional representation and Haar measures, and asserts nonvanishing of a matrix-valued period. It does not assert that the numerical period equals one. Thus a finite admissibility determinant and a period value must remain distinct.

\paragraph{STEP7-02 (positive rank-one period operators).}
The actual normalized factorization in
\path{Development/AbelianVarietyTheory/Definitions/Jets/SumIoPrAlt.lean},
\texttt{IsBeL}, lines 135--139, says, in our conjugate-linear-first convention,
\[
 \operatorname{be}(x,v)\eta=\langle l(v),\eta\rangle l(x)
 \quad(x\in R.\mathrm{sm},\ v\in R.\mathrm{TK}).
\]
Consequently, if this factorization holds and $v$ belongs to both indicated domains, then $P=\operatorname{be}(v,v)$ is positive of rank at most one, with
\[
 \Tr P=\|l(v)\|^2,\qquad \det(I-zP)=1-z\|l(v)\|^2.
\]
\begin{proof}
Put $a=l(v)$, so $P\eta=\langle a,\eta\rangle a$. Cauchy--Schwarz gives $\|P\eta\|\le\|a\|^2\|\eta\|$. Direct calculation gives $\langle\eta,P\eta\rangle=|\langle a,\eta\rangle|^2\ge0$ and $\langle\eta,P\nu\rangle=\langle P\eta,\nu\rangle$. Thus $P$ is bounded and selfadjoint. If $a=0$, all assertions are immediate. Otherwise $a/\|a\|$ is an eigenvector of eigenvalue $\|a\|^2$, and its orthogonal complement is killed. This proves the rank, trace-class, trace and determinant assertions.
\end{proof}
The complete signature of
\path{Development/AbelianVarietyTheory/Proofs/FlagMeasures/BesselPeriodPositivityStatement.lean}, lines 11--12, takes the Bessel maps and these factorization equalities as hypotheses, together with its representation, gap and measure data. The elementary implication just proved does not remove those hypotheses. It is genuine source-based operator positivity once the factorization is supplied, but its determinant has degree one and depends on the chosen vector and period map.

There is also an actual finite scalar relation \texttt{IsCoef} in
\path{Development/AbelianVarietyTheory/Definitions/Jets/SideFam.lean}, lines 30--32:
\[
 F(g)=\sum_i\langle x_i,\sigma(g)y_i\rangle X_i,
 \qquad s=\sum_i\langle l(x_i),X_i l'(y_i)\rangle.
\]
It carries a specified finite matrix-coefficient expansion, not a proved identification with the timeline's $W$. Moreover, the exterior-period decomposition in
\path{Development/AbelianVarietyTheory/Proofs/FlagMeasures/BlockProjectionBesselStatement.lean}, lines 9--10, has coefficients
$\epsilon_I\kappa_{1,I}\overline{\kappa_{0,I}}d_{R_I}$, with $\epsilon_I\in\{\pm1\}$ and formal degrees $d_{R_I}$. The supporting tensor signs are in \texttt{TensH}, \path{SumIoPrAlt.lean}, lines 88--94. These coefficients are not asserted to be $0$ or $1$. This finite exterior-power structure is precisely the kind of algebra recorded by SGA 6, V; its extra analytic period hypotheses and the missing Xi identification must be retained.

\paragraph{DET-08 (the strongest direct determinant link for the toy).}
There \emph{is} a precise finite determinant relation for the path sums. With $P=aa^*$,
\[
 S_k(h)=\Tr(A^{k-1}P),\qquad
 \sum_{k\ge1}S_k(h)z^{k-1}
 =-\left.\partial_u\log\det(I-zA-uP)\right|_{u=0}.
\]
\begin{proof}
Expanding the finite trace gives $\Tr(A^{k-1}aa^*)=a^*A^{k-1}a$. The Neumann series sums it to $a^*(I-zA)^{-1}a$. For invertible $M$, multilinearity of columns gives
$\det(M-uaa^*)=\det M(1-u a^*M^{-1}a)$: terms selecting two rank-one columns vanish because those columns are proportional. Set $M=I-zA$ and differentiate the local logarithm at $u=0$. This proves the sign and the full identity near $(z,u)=(0,0)$.
\end{proof}
This constructive bridge uses an insertion $P$ and a second parameter. It neither replaces path sums by $\Tr A^k$ nor makes $A$ positive, and its finite-dimensional generating function is rational.

\subsection{DET-09: the actual examined RH endpoint is scalar}
An adaptive source trace goes beyond searching for the name $W$. The exported theorem in \path{Supplements/Kernels/MainTheorems.lean}, declarations \texttt{allChi} and \texttt{rhOfAll}, uses the character nonvanishing theorem and \texttt{Formalism.VW.rh\_of\_every}. Following the immediate calls reaches \path{Development/AutomorphicArithmetic/CharacterBoxes/ZeroFreeAllOfTopG.lean}, \texttt{zeroFreeAll\_of\_topG}, lines 26--34, which supplies a scalar bound and \texttt{formB} to \texttt{zeroFreeOn\_of\_tiedWin\_cyc}. The precise path and all intermediate declarations are recorded in the evidence note. The mathematical types at the endpoint are explicit.

The source's smoothed von Mangoldt function is
\[
 G_\zeta(y)=\sum_{1\le n\le\lfloor y\rfloor}\Lambda(n)(y/n-1)
 -y\log y+(1+\gamma_E)y,
\]
from \path{Formalism/StatementOC.lean}, lines 11--13. This directly matches step 5 of the note. A level $\mathcal L$ defines a finite positive real array $w_\psi$ and the scalar
\[
 G_{\mathcal L}(y)=G_\zeta(y)+\sum_{\psi\in\mathcal L.\mathrm{Nt}}w_\psi\Re g_{\psi.\chi}(y).
\]
For nontrivial primitive $\chi$, the function $g_\chi$ is the smoothed character von Mangoldt sum with the correction $yL'(1,\chi)/L(1,\chi)$. The exact definitions are \path{Formalism/SharpRestriction/RamificationFrobenius/TotientGaloisClasses/Analytic.lean}, lines 9--24, and \path{ExplicitFormulaErrorOfHostConditions.lean} in the same directory, lines 32--41. The weights equal either $1$ or $(3\mathcal L.o+1)^{-1}$ with $\mathcal L.o\in\mathbb N$. Their positivity does not assert $G_{\mathcal L}\ge0$.

\path{Development/AutomorphicArithmetic/PrimitiveCosets/CoreGm.lean}, \texttt{core\_gm}, lines 12--15 and 55--74, has a one-sided real conclusion
\[
 G_{\mathcal L}(x)\le x^{1/2+(e_D+e_L+7)\varepsilon}
\]
under its specified arithmetic data. In the family version, \path{Formalism/SharpRestriction/RamificationFrobenius/TotientGaloisClasses/TopWin.lean}, lines 23--32 and 68--93, the exponent approaches $1/2$ as $\varepsilon\downarrow0$, with bounds on power windows $[Y,Y^\lambda]$, $\lambda>1$, at arbitrarily large $Y$. The definitions of \texttt{HasWindows}, \texttt{FormBAt} and \texttt{FormB} in \path{Window.lean} in the same directory, lines 7--8 and 30--33, state the scalar analytic bridge: such $x^\theta$ bounds imply nonvanishing for $\Re s>\theta$ in the retained character family. Choosing $1/2<\theta<\min(\Re s,1)$ excludes a putative right-half-plane zero; functional-equation reflection then gives the critical line.

The implementation of \texttt{formB} is \path{Development/AutomorphicArithmetic/PrimitiveCosets/Imports.lean}, lines 8--19. It invokes \texttt{zeroFree\_of\_rpow\_window} in \path{Development/AutomorphicArithmetic/DedekindLFunctions/LandauTheorems/ZeroFreeOfRpowWindow.lean}, lines 11--14, and \texttt{lemmaL} in the adjacent \path{lemmaL.lean}, lines 13--22. The relevant \texttt{PoleData}, defined in \path{Development/AutomorphicArithmetic/DedekindLFunctions/LandauZeroFreeCriterion.lean}, lines 15--23, is data for a scalar meromorphic Laplace transform
\[
 \Phi(s)=\int_0^\infty G_F(e^u)e^{-su}\,du\qquad(\Re s>2),
\]
where $F$ is a finite weighted character family. It is not an operator trace.

\paragraph{Checked residue calculation.}
The adjacent \path{LandauTheorems/poleData.lean}, lines 137--144, 250--261 and 344--383, gives
\[
 c_\rho=\frac{S_\rho}{\rho(1-\rho)},\qquad
 S_\rho=\sum_i\frac{\alpha_i}{2}\bigl(m_i(\rho)+m'_i(\rho)\bigr)>0,
\]
where the $\alpha_i$ are positive real character weights and the multiplicities are nonnegative integers. The membership of $\rho$ in the pole set provides at least one positive multiplicity, so the finite sum $S_\rho$ is strictly positive. For $\rho\ne0,1$ this proves $c_\rho\ne0$. For a real pole $0<\rho<1$, it also proves $c_\rho>0$. For a general nonreal pole it proves no real-order positivity. These deductions follow directly from the displayed finite sum and denominator; they do not re-prove the analytic Landau theorem.

This is a source-supported description of the final analytic mechanism: arithmetic estimates give scalar window bounds, a scalar Laplace-transform argument excludes zeros, and reflection finishes RH. It is not a claim that every upstream proof has been read. The shared factor $1/[\rho(1-\rho)]$ is exactly the scalar transform of the designated $B$ used in \fid{B-03}, since $\rho(1-\rho)=\gamma^2+1/4$. The residue calculation alone does not identify an operator representation. The further calculation in \fid{DET-11} does give an exact representation of the scalar zeta endpoint by a resolvent trace of this zero-based $C$, with all completion terms retained. The justified implication to the specified $B$ is still
\[
 \text{cited RH}\ \Longrightarrow\ \gamma_j\in\R\ \Longrightarrow\ B\ge0.
\]
Fable's reference to Landau therefore has a correct scalar counterpart in the actual endpoint. Its proposed equality of Fourier operators and Fredholm trace moments does not follow, and was explicitly refuted above. An exact historical declaration for the note's unqualified step-4 $W(1_S)$ remains unmatched; this no longer obstructs identifying the mathematical type of the examined analytic endpoint.

\subsection{DET-11: an exact resolvent bridge for the scalar prime-sum endpoint}
There is a stronger affirmative comparison than the common residue denominator. It concerns the scalar endpoint itself, rather than an identification of the upstream finite matrices. Let $G_\zeta$ be exactly the function in \fid{DET-09}, and put
\[
 \Phi_\zeta(s)=\int_0^\infty G_\zeta(e^u)e^{-su}\,du,
 \qquad v=s(1-s),\qquad C=B(I+B/4)^{-1}.
\]
Write $\psi=\Gamma'/\Gamma$ and
\[
 A_\infty(s)=\frac1s+\frac1{s-1}-\frac12\log\pi+
 \frac12\psi(s/2).
\]
The following identities are unconditional, with the zero-based $B$ of \fid{B-01}.
\begin{proposition}
The trace-class determinant $D_C(v)=\det(I-vC)$ satisfies, for every $s\in\C$,
\begin{equation}\label{eq:Cdetcompletion}
 D_C(s(1-s))=2\xi(s).
\end{equation}
Where $\xi(s)\ne0$,
\begin{equation}\label{eq:Cresolvent}
 \frac{\xi'(s)}{\xi(s)}=(2s-1)\Tr\bigl(C(I-s(1-s)C)^{-1}\bigr).
\end{equation}
In particular, for $\Re s>1$ the actual scalar transform is
\begin{align}\label{eq:Phiresolvent}
 \Phi_\zeta(s)
 &=\frac{A_\infty(s)-(2s-1)\Tr\bigl(C(I-vC)^{-1}\bigr)}{s(s-1)}\notag\\
 &\hspace{7mm}-\frac1{(s-1)^2}+\frac{1+\gamma_E}{s-1}.
\end{align}
Thus its nontrivial-zero contribution is an ordinary trace of a resolvent of a fixed trace-class operator after explicit completion terms. This representation uses the zeros as input.
\end{proposition}
\begin{proof}
The diagonal entries of $I+B/4$ are nonzero: a zero would give $\gamma^2=-1/4$, hence a zero of $\xi$ at $0$ or $1$, whereas $\xi(0)=\xi(1)=1/2$. They tend to one, so the inverse is bounded. The entries of $C$ are therefore
\[
 c_j=(\gamma_j^2+1/4)^{-1}=[\rho_j(1-\rho_j)]^{-1},
 \qquad \sum_j|c_j|<\infty.
\]
Absolute locally uniform convergence of the products gives
\[
 D_C(v)=\frac{\prod_j(1-(v-1/4)\gamma_j^{-2})}
 {\prod_j(1+\gamma_j^{-2}/4)}
 =\frac{G(v-1/4)}{G(-1/4)}.
\]
Since $v-1/4=-(s-1/2)^2$, evenness of $\Xi$ gives
$G(v-1/4)=\xi(s)/\xi(1/2)$ and
$G(-1/4)=\xi(0)/\xi(1/2)$. This proves \eqref{eq:Cdetcompletion}, including the normalization factor two and with no square-root branch choice.

On any compact set avoiding the zeros of $D_C$, the sum
$\sum_j c_j/(1-vc_j)$ converges uniformly and absolutely. Indeed, the tail denominators are uniformly bounded away from zero because $c_j\to0$, and the finitely many remaining denominators are nonzero there. Differentiating the locally convergent logarithmic products consequently gives
\[
 \frac{D_C'(v)}{D_C(v)}=-\sum_j\frac{c_j}{1-vc_j}
 =-\Tr\bigl(C(I-vC)^{-1}\bigr).
\]
The inverse is bounded by the same diagonal argument, and its product with $C$ is trace class. The chain rule and $v'(s)=1-2s$ prove \eqref{eq:Cresolvent}. At $s=1/2$ no division by $2s-1$ is made; both sides of that identity vanish.

The source already computes the Mellin transform in
\path{Development/AutomorphicArithmetic/DedekindLFunctions/LandauTheorems/mellinGchiModOne.lean},
\texttt{mellin\_Gchi\_modOne}, lines 115--118. Its change from $y$ to $u=\log y$ is
\texttt{poleData\_change}, lines 114--135 of the adjacent \path{poleData.lean}.
For the modulus-one character its displayed formula is algebraically
\[
 \Phi_\zeta(s)=-\frac{\zeta'(s)/\zeta(s)}{s(s-1)}
 -\frac1{(s-1)^2}+\frac{1+\gamma_E}{s-1},\qquad \Re s>1.
\]
The modulus-one definition is \path{CharFamily.lean}, lines 15--22 in the parent directory; its ramp equals the sum in \path{Formalism/StatementOC.lean}, lines 11--13 for $y\ge1$, because terms beyond $\lfloor y\rfloor$ vanish and $\Lambda(0)=\Lambda(1)=0$.
Logarithmic differentiation of the defining completion gives
$\xi'/\xi=A_\infty+\zeta'/\zeta$. Substitute \eqref{eq:Cresolvent} to obtain \eqref{eq:Phiresolvent}. The convergence and transform statement are those of the cited source on $\Re s>1$; its formula agrees with direct integration of each ramp term, whose Mellin transform is $n^{-s}/[s(s-1)]$.
\end{proof}

This finding narrows the comparison: the source's scalar endpoint does have an exact Fredholm-resolvent representation, so an assertion that the connection stops at a similar denominator would be false. It does not identify any finite Gram, Capelli or period matrix with $C$ or $B$. It also does not make the source's scalar inequalities into a direct proof of selfadjointness of an arithmetic operator: the representation was built from the same zero data as \fid{B-01}. The passage from $C$ back to $B$ is the scalar spectral transform $B=C(I-C/4)^{-1}$, with bounded inverse because $c_j\ne4$ and $c_j\to0$. Geometrically, the determinant is still the analytic exterior-power construction of cote 1 and SGA 6, V, rather than a newly constructed cohomology over $\operatorname{Spec}\Z$.

\subsection{STEP6-01--03: an exact candidate for the two-indicator Weil test}
The timeline does not define its $W$. We can nevertheless give a precise candidate normalization, identify its arithmetic cross term, and prove the stated half-plane shape. None of these calculations attributes this normalization to the historical determinant construction.

Put $g=\mathbf1_{(0,1]}$, $g_X=\mathbf1_{(0,X^{-1}]}$ and $f^\pm_{X,\eta}=g\pm X^\eta g_X$, for $X\ge1$ and $\eta\in\R$. On these real functions define
\[
 \M f(s)=\int_0^\infty f(x)x^s\,\frac{dx}{x}\quad(\Re s>0),
 \qquad W_{\rm std}(f)=\sum_\rho\M f(\rho)\M f(1-\rho),
\]
with all nontrivial zeros and their multiplicities. This is a quadratic form in $f$. If $W$ instead denotes the linear zero functional, this expression applies it to $f*_\times f^\sharp$, where $f^\sharp(x)=x^{-1}\overline{f(x^{-1})}$: its Mellin transform is $\M f(s)\overline{\M f(1-\bar s)}$. Fubini is valid for these functions on $0<\Re s<1$. For example, $g*_\times g^\sharp(x)=\min(1,x^{-1})$, not $g(x)$.

\paragraph{STEP6-01: the trace and its positive product.}
Use the paired-zero Hilbert space of \fid{B-01}, put
\[
 Ae_j=\gamma_je_j,\qquad
 D(A)=\{v:\sum_j|\gamma_j|^2|v_j|^2<\infty\},\qquad
 C=(A^2+I/4)^{-1}=B(I+B/4)^{-1}.
\]
Here $A$ is the maximal closed normal diagonal operator. For $L=\log X$ and $a=X^{\eta-1/2}$,
\begin{equation}\label{eq:WeilIndicatorTrace}
 W_{\rm std}(f^\pm_{X,\eta})
 =2\Tr\!\left[C\bigl((1+a^2)I\pm2a\cos(AL)\bigr)\right],
 \qquad W_{\rm std}(g)=2\Tr C=2\lambda_1.
\end{equation}
Under RH, $A$ is selfadjoint, $U=e^{iAL}$ is unitary and, writing $F_\pm=I\pm aU$,
\begin{equation}\label{eq:WeilIndicatorPositive}
 W_{\rm std}(f^\pm_{X,\eta})
 =2\Tr(CF_\pm^*F_\pm)
 =2\|F_\pm C^{1/2}\|_{\mathcal S_2}^{2}\ge0.
\end{equation}
\begin{proof}
Direct absolutely convergent integration gives $\M f^\pm(s)=(1\pm X^{\eta-s})/s$ for $\Re s>0$. If $\rho=\beta+it$ is a nontrivial zero, the numerator product is bounded by $(1+X^\eta)^2$ and $|\rho(1-\rho)|\ge t^2$. The zero count $N(T)=O(T\log T)$ gives
$\sum_\rho|\rho(1-\rho)|^{-1}<\infty$: the dyadic annulus $2^m\le|t|<2^{m+1}$ contributes $O((m+1)2^{-m})$, and the finitely many low zeros have nonzero denominators. Thus the quadratic sum is absolutely convergent, locally uniformly in $X,\eta$, and can be paired or conjugated.

The entries $c_j=(\gamma_j^2+1/4)^{-1}=[\rho_j(1-\rho_j)]^{-1}$ define a trace-class $C$. The sequences $c_j$ and $\gamma_j^2c_j$ are bounded: the denominators never vanish and $|\gamma_j|\to\infty$. Hence $C$ maps the Hilbert space into $D(A^2)$ and is the two-sided inverse of $A^2+I/4$ on the appropriate domains. Since $|\Im\gamma_j|<1/2$, the diagonal operators $e^{\pm iAL}$ and $\cos(AL)$ are bounded, with entry bounds $X^{1/2}$ and $\cosh(L/2)$. Their products with $C$ are trace class.

For either sign, expansion gives
\[
 \M f^\pm(\rho)\M f^\pm(1-\rho)
 =\frac{1+X^{2\eta-1}\pm X^{\eta-\rho}\pm X^{\eta-1+\rho}}
 {\rho(1-\rho)}
 =\frac{1+a^2\pm2a\cos(\gamma L)}{\gamma^2+1/4}.
\]
This is even in $\gamma$, so every reflection pair contributes twice the representative's diagonal entry. It proves the first identity in \eqref{eq:WeilIndicatorTrace}. The base test gives $2\Tr C$, and \fid{B-03} gives $\Tr C=\lambda_1$.

Under RH, $C\ge0$ and all these diagonal operators commute. Therefore
$F_\pm^*F_\pm=(1+a^2)I\pm a(U+U^*)=(1+a^2)I\pm2a\cos(AL)$.
The operator $C^{1/2}$ is Hilbert--Schmidt. Taking the trace of
$(F_\pm C^{1/2})^*(F_\pm C^{1/2})$ proves \eqref{eq:WeilIndicatorPositive}. Off RH, the exact trace remains valid, but $U^{-1}$ need not be $U^*$ and $C$ need not be positive; the positive-product step is not asserted there.
\end{proof}

\paragraph{STEP6-02: the actual smoothed sum is a weighted cosine trace.}
Let $G_\zeta(X)$ be the exact source scalar of \fid{DET-09}, namely
$\sum_{n\le X}\Lambda(n)(X/n-1)-X\log X+(1+\gamma_E)X$. Set
\[
 R(X)=\sum_\rho\frac{X^\rho}{\rho(1-\rho)},\qquad
 b(X)=\log(2\pi)-\sum_{n=1}^\infty\frac{X^{-2n}}{2n(2n+1)}.
\]
Unconditionally, for every $X\ge1$,
\begin{equation}\label{eq:GzCosine}
 G_\zeta(X)=b(X)+R(X)
 =b(X)+2\sqrt X\,\Tr\bigl(C\cos(A\log X)\bigr).
\end{equation}
All sums converge absolutely and locally uniformly. The correction satisfies
\[
 b(1)=\log(4\pi)-1,\qquad
 \log(4\pi)-1\le b(X)<\log(2\pi).
\]
Under RH, \eqref{eq:GzCosine} gives
$|G_\zeta(X)-b(X)|\le2\lambda_1\sqrt X$.
\begin{proof}
The primary explicit-formula input is the \emph{corrected} Lemma 2.1 of Ramar\'e's corrigendum \cite[pp.~1--2]{ramare}; the original article's Lemma 2.2 has sign errors. With $\psi(X)=\sum_{n\le X}\Lambda(n)$ and $\widetilde\psi(X)=\sum_{n\le X}\Lambda(n)/n$, the corrected formula is
\[
 \widetilde\psi(X)=\log X-\gamma_E+\frac{\psi(X)-X}{X}
 -\sum_\rho\frac{X^{\rho-1}}{\rho(\rho-1)}+\frac{b(X)}X.
\]
Multiply by $X$, subtract $\psi(X)$ and then subtract
$X\log X-(1+\gamma_E)X$. This proves $G_\zeta=R+b$, because the source ramp is $X\widetilde\psi-\psi$. The zero sum is absolutely convergent by $|X^\rho|\le X$ and the reciprocal-denominator estimate above. Pairing $\rho$ with $1-\rho$ gives
$X^\rho+X^{1-\rho}=2\sqrt X\cos(\gamma\log X)$ and proves the trace identity.

The corrigendum also gives
\[
 b(X)=X\int_X^\infty
 \left(\log(2\pi)+\tfrac12\log(1-t^{-2})\right)\frac{dt}{t^2}.
\]
Expand $-\log(1-t^{-2})=\sum_{n\ge1}t^{-2n}/n$ and apply Tonelli to the nonnegative series. This gives exactly the displayed series for $b$, including $X=1$, where the logarithmic singularity is integrable. The identity
$1/[2n(2n+1)]=1/(2n)-1/(2n+1)$ gives the sum $1-\log2$ at $X=1$. Positivity and monotonicity of its terms prove the bounds and continuity. For $X>1$ the equivalent elementary expression is
\[
 b(X)=\log(2\pi)+\tfrac12\log(1-X^{-2})
 +\tfrac12\left(X\log\frac{X+1}{X-1}-2\right).
\]
At prime-power endpoints the jumps of $X\widetilde\psi-\psi$ cancel, so no half-weight convention remains in $G_\zeta$. Setting $X=1$ checks
$2\lambda_1=G_\zeta(1)-b(1)=2+\gamma_E-\log(4\pi)$.
Under RH, $|\cos(\gamma_j\log X)|\le1$ and $c_j>0$, giving the asserted bound.
\end{proof}
This derivation cites the classical analytic explicit formula; it independently proves its algebraic translation, convergence, endpoint handling and trace representation. It does not construct $A$ arithmetically from a period determinant.

\paragraph{STEP6-03: exact one-sided bounds and the half-plane implication.}
The two reflected cross sums coincide, so
\begin{align}
 W_{\rm std}(f^\pm_{X,\eta})
 &=2\lambda_1(1+X^{2\eta-1})
   \pm2X^{\eta-1}\bigl(G_\zeta(X)-b(X)\bigr),\label{eq:WeilGz}\\
 W_{\rm std}(f^+_{X,\eta})\ge0
 &\Longleftrightarrow G_\zeta(X)\ge b(X)-\lambda_1(X^{1-\eta}+X^\eta),\label{eq:WeilLower}\\
 W_{\rm std}(f^-_{X,\eta})\ge0
 &\Longleftrightarrow G_\zeta(X)\le b(X)+\lambda_1(X^{1-\eta}+X^\eta).\label{eq:WeilUpper}
\end{align}
Thus the plus sign gives a lower bound, while the minus sign gives an upper bound. At $X=1$ the values are $8\lambda_1$ and $0$, consistent with $f^+=2g$ and $f^-=0$.

\begin{lemma}
If $0<\theta<1$ and either $G_\zeta(X)\ge-KX^\theta$ or $G_\zeta(X)\le KX^\theta$ for some $K\ge0$ and all sufficiently large $X$, then $\zeta$ has no zero with real part greater than $\theta$.
\end{lemma}
\begin{proof}
The scalar Mellin identity established in \fid{DET-11} is
\[
 \Phi_\zeta(s)=-\frac{\zeta'(s)/\zeta(s)}{s(s-1)}
 -\frac1{(s-1)^2}+\frac{1+\gamma_E}{s-1}\qquad(\Re s>1).
\]
It continues meromorphically to $\Re s>0$. At $s=1$, putting $h=s-1$, the expansion
$-\zeta'/\zeta=h^{-1}-\gamma_E+O(h)$ makes the first term
$h^{-2}-(1+\gamma_E)h^{-1}+O(1)$, so the polar terms cancel. At a nontrivial zero $\rho$ of multiplicity $m$, its residue is $m/[\rho(1-\rho)]\ne0$. There are no real zeros on $(0,1)$, since the positive alternating eta series equals $(1-2^{1-s})\zeta(s)$ there; for $s>1$ use the Euler product. Thus $\Phi_\zeta$ is analytic near every positive real point.

For a lower bound choose $u_0\ge0$ beyond its threshold and define $h(u)=0$ before $u_0$ and $h(u)=G_\zeta(e^u)+Ke^{\theta u}\ge0$ afterwards. The crude bound $G_\zeta(X)=O(X(1+\log X)^2)$ follows from $\Lambda(n)\le\log n$. Hence its Laplace transform $H$ converges for $\Re s>1$, where
\[
 H(s)=\Phi_\zeta(s)-\int_0^{u_0}G_\zeta(e^u)e^{-su}\,du
       +\frac{Ke^{-(s-\theta)u_0}}{s-\theta}.
\]
The compact integral is entire. This gives meromorphic continuation of $H$ to $\Re s>\theta$, analytic near every real point there, with nonzero poles at any zeta zeros there.

For clarity, the required positive-Laplace fact has a short proof. If a nonnegative locally integrable function has finite abscissa of convergence $\sigma_c$, its Laplace transform is singular at the real point $\sigma_c$. Otherwise let a disc of analytic continuation there have radius $\delta>0$. Set $\sigma_0=\sigma_c+\delta/4$ and $r=\delta/2$. The Taylor disc about $\sigma_0$ has radius at least $3\delta/4$, whereas $\sigma_0-r<\sigma_c$. Differentiation under the integral at $\sigma_0$ is justified by convergence at a slightly smaller point still above $\sigma_c$. Its Taylor series towards the left has nonnegative terms, and Tonelli gives
\[
 \sum_{n\ge0}\frac{r^n}{n!}
  \int_0^\infty u^nh(u)e^{-\sigma_0u}\,du
 =\int_0^\infty h(u)e^{-(\sigma_0-r)u}\,du<\infty,
\]
contradicting the definition of $\sigma_c$. The case $\sigma_c=-\infty$ already converges everywhere.

Here $\sigma_c\le1$. If $\sigma_c>\theta$, the displayed continuation is analytic through that real point, contradicting this fact. Thus $\sigma_c\le\theta$, so $H$ is holomorphic on $\Re s>\theta$. Meromorphic uniqueness forces $\Phi_\zeta$ to have no pole there, excluding all the indicated zeros. For an upper bound use $h(u)=Ke^{\theta u}-G_\zeta(e^u)$ beyond $u_0$; the same proof applies to $-\Phi_\zeta$ and its nonzero residues.
\end{proof}

For $0<\eta<1$, eventual nonnegativity of either specified test therefore gives zero-freeness in $\Re s>\max(\eta,1-\eta)$, by \eqref{eq:WeilLower}--\eqref{eq:WeilUpper} and boundedness of $b$. In particular $0<\eta\le1/2$ gives exactly $\Re s>1-\eta$. At $\eta=1/2$ eventual positivity of the plus-test family is RH-equivalent: the forward direction uses the lemma and reflection, and the reverse is \eqref{eq:WeilIndicatorPositive}. This proves an equivalence for the defined form, not the unprovided historical identification.

The approximate amplitude is also transparent. For any real $q(X)>0$, the same expansion gives
$W_{\rm std}(g+qg_X)=2\lambda_1(1+q^2/X)+2qR(X)/X$.
Its positivity is exactly $G_\zeta(X)\ge b(X)-\lambda_1(X/q+q)$.
If $c_1X^\eta\le q(X)\le c_2X^\eta$ eventually, for positive constants $c_1,c_2$, the same exponent follows, using $|\lambda_1|$ in the coarse constant.

The remaining source gate is concrete: specify the linear or quadratic convention for $W$, the test correspondence and normalization, and the passage from the finite period/determinant construction to this zero form. The source scalar is exactly \path{Formalism/StatementOC.lean}, lines 11--13. The actual final path in \fid{DET-09} bounds a character sum with positive weights \emph{from above on power windows}; the present lemma proves a zeta-only criterion for a global eventual bound. The source also has lower bounds on differences $G_\zeta-G_\chi$, which are a different assertion. Character averaging and window quantifiers remain additional parts of any correspondence with that path. The positive product here uses the stipulated RH; the converse lemma does not prove its own positivity premise. These findings shrink the mathematical interpretation of step 6 while leaving its historical attribution and the arithmetic construction of the operator open.

\section{TOY-01--02: the estimate, its domain and the invalid sign inference}
\begin{theorem}
Let $q$ be a prime power, $t:\F_q\to[0,\infty)$, and $k\ge2$. In counting measure put
\[
 Q_k(x)=\sum_{i=1}^{k-1}x_ix_{i+1},\quad
 N_k=\sum_{Q_k(x)=0}\prod_{i=1}^kt(x_i),\quad M=\sum_x t(x).
\]
Then $qN_k=M^k+R_k$, with
\begin{equation}\label{eq:toybound}
 |R_k|\le(q-1)q^{(k-1)/2}\|t\|_2^2\|t\|_\infty^{k-2}.
\end{equation}
The constant is attained for $t\equiv1$ and odd $k\ge3$.
\end{theorem}
\begin{proof}
Choose a nontrivial additive character $\psi$. Orthogonality gives
$q1_{Q=0}=\sum_{h\in\F_q}\psi(hQ)$. The zero mode contributes $M^k$. For $h\ne0$ the other mode is
\[
 S_k(h)=\langle t,K_hD_tK_h\cdots D_tK_ht\rangle,
\]
with $k-1$ kernels and $k-2$ multipliers. Since
$(K_hK_h^*)_{x,x'}=\sum_y\psi(h(x-x')y)=q\delta_{x,x'}$,
$\|K_h\|=\sqrt q$, and $\|D_t\|=\|t\|_\infty$. Cauchy--Schwarz and the operator-norm inequality give the summand bound in \eqref{eq:toybound}; summing over $q-1$ modes proves it. At $k=2$ the multiplier product is empty. The zero weight is handled directly.

For $t\equiv1$, $K_h^2=qJ$ with $J1=1$. If $k=2m+1$, then $S_k=1^TK_h^{2m}1=q^{m+1}$, independent of $h\ne0$. This is equality in the bound, including its final triangle inequality.
\end{proof}
The original Fable statement specifies $q$ prime, so it is correct with the added $k\ge2$ convention. The prompt's shorter phrase ``on $\Z/q$'' needs the following distinction.

\begin{proposition}
For arbitrary integer $q\ge2$ on the cyclic ring $\Z/q$, a valid replacement is
\begin{equation}\label{eq:composite}
 |R_k|\le q^{(k-1)/2}\|t\|_2^2\|t\|_\infty^{k-2}
 \sum_{h=1}^{q-1}\gcd(h,q)^{(k-1)/2}.
\end{equation}
The prime-form bound fails at $q=4$, $t\equiv1$, $k=3$. The displayed bound also fails at $k=1$, even for prime $q$.
\end{proposition}
\begin{proof}
Set $d_h=\gcd(h,q)$. The entries of $K_hK_h^*$ equal $q$ when $x-x'$ is divisible by $q/d_h$ and zero otherwise. After a permutation this is a sum of $d_h\times d_h$ all-ones blocks multiplied by $q$. Hence $\|K_h\|=\sqrt{qd_h}$, and the preceding proof gives \eqref{eq:composite}.

For $q=4$, $Q_3=x_2(x_1+x_3)$. The values $x_2=0,1,2,3$ allow respectively $16,4,8,4$ choices, so $N_3=32$, $R_3=4\cdot32-4^3=64$. The claimed prime-form right side is $3\cdot4\cdot4=48$. The corrected bound is $16(1+2+1)=64$. At $k=1$, $N_1=M$, $R_1=(q-1)M$. For $q=2$, $t=(2,1)$ this is $3$, but the claimed expression is $5/2$.
\end{proof}
The exact computation script checks 144 finite prime-field cases and the two counterexamples. It corroborates these proofs and is not a proof of their universal statements.

\paragraph{TOY-02 (scalar signs do not imply spectral positivity).}
Even the stronger condition $\Tr T^n>0$ for every $n\ge1$ fails to imply $T\ge0$: take $T=\diag(1,-1/2)$. Then $\Tr T^n=1+(-1/2)^n>0$, but $\det(1-wT)=(1-w)(1+w/2)$ has a negative zero. Thus the positive-axis singularity statement for a power series with nonnegative coefficients (often called Pringsheim's theorem in this setting) does not supply Fable's proposed conclusion. In a cohomological no-cancellation argument, the degrees, weight bounds and factorization are substantive extra hypotheses.

\paragraph{Geometric location.}
For $t\equiv1$, $N_k$ counts points of the affine $\F_q$-scheme $Q_k=0$. If $t$ is realized as the trace function of a constructible $\ell$-adic sheaf $\mathcal T$, the twisted sums have the usual interpretation using $\mathcal T^{\boxtimes k}$ and the Artin--Schreier character sheaf pulled back by $hQ_k$. The trace formalism is SGA 5, III, Corollaries 4.7--4.8; its curve specialization is III B, Corollary 1.3.1 \cite{sga5}. An arbitrary real nonnegative function $t$ is not automatically such an arithmetic sheaf. The proof above is finite Fourier analysis and does not construct a motive over $\Spec\Z$, a crystal, or divided powers. Exterior powers in the determinant formula must not be confused with divided-power multiplicity in the other prompts.

\section{Masullo: a conditional trace-class application and a different arithmetic identity}
\subsection{The Hilbert quotient and Theorem 10.3.2.1}
Let $F$ be a number field, $C_F=\A_F^\times/F^\times$, and $C_F^1=\ker|\cdot|_F$. For $I=[-\alpha,\alpha]$, $\alpha\ge0$, $a\in\R$ and real $s\ge0$, the draft uses a Hilbert completion with norm
\begin{equation}\label{eq:masnorm}
 \|f\|_{I,s}^2=\int_{C_F}|f(g)|^2
 \max\{|g|_F^{-2\alpha},|g|_F^{2\alpha}\}
 (1+|\log|g|_F|^2)^{s/2}\,d\mu(g).
\end{equation}
Put $\mathcal S_0(\A_F)=\{f\in\mathcal S(\A_F):f(0)=\widehat f(0)=0\}$ and $\Theta_f(g)=\sum_{\xi\in F^\times}f(\xi g)$. Writing $V^s_{a,I}$ for the closure in $W_I^s(C_F)$ of $|\cdot|_F^a$ times the theta-image, the quotient is
\[
 E^s_{F,a,I}=W_I^s(C_F)/V^s_{a,I}.
\]
The draft denotes it $\mathcal H^s_{F,a,I}$. These definitions are in \cite[\S1.1.3; Definition 10.1.1.1; (10.2.2.1)]{masullo}; the norm and dual norm are (9.2.2.34), (9.2.2.40)--(9.2.2.41), p.~417. The notation for continuous homology of the theta-image is not \emph{\'etale} $H^0$ of an arithmetic scheme.

The test-function action descends from convolution,
\[
 A(\phi)f(x)=\int_{C_F}\phi(g)f(g^{-1}x)\,d\mu(g),\qquad \phi\in\cD(C_F).
\]
\textbf{Cited draft assertion MAS-SRC-01.} Theorem 10.3.2.1, ``Spectral Nuclearity II,'' pp.~453--455, asserts for $s>1$ convergence of finite-rank expansions in every Schatten class $\mathcal S_p$, $p\ge1$, and specifically
\[
 A(\phi)\in\Sone(E^s_{F,a,I}),\qquad
 \Tr A(\phi)=\sum_{\chi,t,b}N_{I,s}(\rho_{\chi,t,b})
 \M\phi(\chi|\cdot|_F^{b+it}).
\]
Here $\M\phi(c)=\int\phi(g)c(g)d\mu(g)$; $\rho_{\chi,t,b}=\chi|\cdot|^{a+b+it}$ is a zero in the closed strip $a+I$. Interior multiplicity is full; the printed boundary cap is discussed below. The exact displays are (10.3.2.1)--(10.3.2.4); the character-isotype formula is (10.3.2.5). The separate central-line version is Theorem 10.3.1.1. The spectral quotient basis of Theorem 10.2.2.1 is described as a Markushevich basis; the orbital Schauder bases belong to Chapter 7. The two basis constructions should not be conflated.

\paragraph{MAS-01 (conditional application of cote 1).}
Conditional on the asserted trace-class theorem, each $A(\phi)$ has an ordinary Hilbert-space Fredholm determinant and satisfies
\begin{align*}
 D_\phi(w)&=\det(1-wA(\phi))=\prod_j(1-w\lambda_j(A(\phi))),\\
 \log D_\phi(w)&=-\sum_{n\ge1}\frac{w^n}{n}\Tr A(\phi)^n
 \quad (|w|\,\|A(\phi)\|<1).
\end{align*}
\begin{proof}
The quotient by the closed subspace is Hilbert, and the cited hypothesis is exactly $\mathcal S_1$. Thus the cote 1 trace/determinant theorem applies without any upgrade from a merely basis-defined locally convex trace. The convolution representation gives $A(\phi)^n=A(\phi^{*n})$. This is an application of the known theorem, not a proof of the draft's analytic estimates.
\end{proof}
This is the affirmative trace-theoretic comparison. The novel arithmetic claim lies in relating that trace to the geometric orbital terms, with the intended independence of arbitrarily narrow positive-width strips.

\paragraph{MAS-02 (not the single determinant of $B$).}
This family cannot equal the nonconstant $G(w)$ for every $\phi$: $A(0)=0$ gives $D_0=1$, and $A(t\phi)=tA(\phi)$ gives $D_{t\phi}(w)=D_\phi(tw)$. A separate transform needs its own spectral theorem. The new \fid{MAS-10} below supplies one such ordinary determinant on the centered zeta quotient, conditional on the draft's full spectral-jet description. This does not identify the whole linear test family with one fixed determinant.

\subsection{Theorem 12.5.1.1: the main arithmetic trace formula}
\textbf{Cited draft assertion MAS-SRC-02.} The theorem entitled ``Trace Formula'' at \cite[Theorem 12.5.1.1, pp.~519--520, (12.5.1.1)--(12.5.1.3)]{masullo} assumes
\[
 \phi\in\cD(C_F),\qquad \phi|_{C_F^1}=0.
\]
With the local Haar measures and principal values of that draft, it states
\begin{align}
 &\sum_{v\text{ real}}\frac12\int_{F_v^\times}'
       \frac{\phi(g_v)|g_v|_v^{1/2}}{|1-g_v|}\,d\mu_v(g_v)
 +\sum_{v\text{ complex}}\frac1{2\pi}\int_{F_v^\times}'
       \frac{\phi(g_v)|g_v|_v^{1/2}}{|1-g_v|}\,d\mu_v(g_v)\notag\\
 &\quad+\sum_{v\text{ finite}}\frac{\log q_v}{\mu_v(\mathcal O_v^\times)}
       \int_{F_v^\times}'\frac{\phi(g_v)|g_v|_v^{1/2}}{|1-g_v|_v}\,d\mu_v(g_v)\notag\\
 &=\M\phi(|\cdot|_F^{1/2})+\M\phi(|\cdot|_F^{-1/2})
   -\sum_{c,I}N(c)\M\phi(c|\cdot|_F^{-1/2}).\label{eq:masmain}
\end{align}
The local element is mapped into the idele class group when evaluating $\phi$. The sum is over zeros of the completed Hecke function in the interior of $\sigma(c)\in1/2+I$; $N(c)$ is the full multiplicity. We retain the source's archimedean absolute-value notation rather than silently replacing its normalization.

The displayed theorem leaves $I$ unquantified. The intended range $I=[-\alpha,\alpha]$, $\alpha>0$ arbitrary, occurs in \S\S1.2.1 and 1.4.5--1.4.8, pp.~13--15 and 29--31. This is an explicit statement-level qualification, not a parameter that can be erased. The geometric calculation uses cutoffs of $P^{\mathrm{geom}}_{\rho,1/2}A(\phi)$ as $\Lambda\to\infty$ with $\varepsilon\Lambda=o(1)$. Definitions 10.4.1.1 and 10.4.1.3, pp.~456--457, and Definition 10.4.3.7, (10.4.3.49)--(10.4.3.50), pp.~473--474, give its projections and the complementary kernel
$k_{\mathrm{geom}}(x,y)=\phi(xy^{-1})-k_{\mathrm{spec}}(x,y)$.
The untruncated convolution/geometric operator need not be nuclear in that setup. A limiting geometric diagonal trace must not be replaced by an ordinary trace of that uncut operator.

\paragraph{MAS-06 (a concrete obstruction to the uncut convolution trace).}
The distinction already occurs on the radial part of the full weighted space. Put
$w(u)=e^{2\alpha|u|}(1+u^2)^{s/2}$ and choose a nonzero $h\in C_c^\infty(\R)$ with $h(0)=0$. Convolution by $h$ on $L^2(\R,w(u)du)$ is not compact, hence is not trace class.
\begin{proof}
The weight ratios $w(u+v)/w(u)$ are uniformly bounded for $v$ in a fixed compact set, so the compactly supported convolution is bounded by the weighted version of Young's inequality (conjugate by $f\mapsto\sqrt w f$ and dominate by a compactly supported convolution kernel).
Choose $f(u)=\overline{h(-u)}$. Then $k=h*f$ is compactly supported and continuous, with $k(0)=\int|h|^2>0$. For $R\to+\infty$, put $f_R(u)=f(u-R)/\|f(\cdot-R)\|_{L^2(w)}$. Uniformly for $v$ in any fixed compact set,
\[
 \frac{w(R+v)}{w(R)}\longrightarrow e^{2\alpha v}.
\]
Therefore
\[
 \|h*f_R\|_{L^2(w)}^2\longrightarrow
 \frac{\int|k(v)|^2e^{2\alpha v}dv}{\int|f(v)|^2e^{2\alpha v}dv}>0.
\]
Take $R_j\to\infty$ far enough apart that the translates of $\operatorname{supp}k$ are disjoint. The images $h*f_{R_j}$ are pairwise orthogonal and have norms bounded below, so they have no convergent subsequence, while each $f_{R_j}$ has norm one. This disproves compactness.
\end{proof}
Already for $F=\Q$, one has $C_\Q\simeq\R_{>0}\times\widehat\Z^\times$, with norm the first coordinate. Indeed a unique rational factor clears the finitely many nonzero finite valuations and makes the real coordinate positive; a rational number preserving all finite units and that sign must be one. This normalization and its inverse are continuous on the open unit charts, proving the topological decomposition. The compact factor can be given Haar mass one, so functions constant along it give exactly the closed radial subspace above, with $u=\log|g|$. The same decomposition is recorded in the draft, \S A.2.2.1, (A.2.2.3), p.~546; the Haar product convention is \S A.3.3.1, p.~561. The test $\phi(g)=h(\log|g|_\Q)$ is compactly supported on $C_\Q$ and satisfies $\phi|_{C_\Q^1}=0$. Thus even an admissible test can give a nonnuclear convolution on the full weighted space. This does not contradict nuclearity on the theta quotient: the quotient is a different operator realization. It also does not by itself identify the precise uncut geometric projection with this radial restriction; that identification is not needed for the obstruction to replacing cutoff traces by a full convolution trace.

\textbf{Cited draft assertion MAS-SRC-03.} Theorem 13.3.1.1, p.~522, is a different formulation: it sums only over $\sigma(c)=1/2$, uses $N_{\mathrm{spec}}(c)=\min\{N(c),[F:\Q]\}$, and imposes the Mellin-jet vanishings (13.3.1.2). It cannot be substituted for \eqref{eq:masmain} with full multiplicity. The printed page has a correction instruction for that statement; Theorem 13.3.1.2 has a title but no mathematical assertion. We cite the general-Hecke GRH consequence as the draft's introductory claim and proposed deduction, conditional on it, not as a fully stated numbered theorem 13.3.1.2 checked in this run.

\subsection{MAS-03--05: what positivity means here}
For centered multiplicative convolution write $\psi^*(g)=\overline{\psi(g^{-1})}$. Then
\begin{equation}\label{eq:squareone}
 (\psi*\psi^*)(1)=\int_{C_F}|\psi(g)|^2d\mu(g).
\end{equation}
\begin{proof}
Insert $g=1$ in the convolution integral. Since $C_F$ is abelian, Haar measure is invariant under the inversion used in the involution. The integrand is exactly $|\psi|^2$. Compact support makes the integral finite; if it is zero, continuity makes $\psi$ identically zero.
\end{proof}
Since $1\in C_F^1$, the class admitted in \eqref{eq:masmain} contains no nonzero convolution square. This proves \fid{MAS-03}: the theorem cannot be directly evaluated on the Weil positivity cone. The draft itself makes this distinction in \S1.4.6, p.~30, saying that positivity is not used and returns only after the proposed GRH deduction.

Under GRH, $\eta_c=c|\cdot|^{-1/2}$ is unitary. In a fixed smooth character component the usual convergent Weil zero functional therefore satisfies
\begin{equation}\label{eq:weilsquares}
 W(\psi*\psi^*)=\sum_cN(c)|\M\psi(\eta_c)|^2\ge0.
\end{equation}
\begin{proof}
Fubini gives multiplicativity of the Mellin transform on convolution. Inversion gives
$\M(\psi^*)(\eta)=\overline{\M\psi(\overline\eta^{-1})}$, which is $\overline{\M\psi(\eta)}$ for unitary $\eta$. The Mellin transform of a smooth compactly supported function decays faster than every power in vertical strips, and the standard fixed-$L$-function zero count is $O(T\log T)$, so the zero sum converges absolutely. Summing the nonnegative terms proves the formula.
\end{proof}
This is \fid{MAS-04}. For $\zeta$ and Dirichlet $L$-functions it is an Alp\"oge consequence. For all Hecke $L$-functions of all number fields it is conditional on the supplied Masullo claim. It is scalar quadratic positivity. Masullo's weighted translation action is explicitly nonunitary (Remark 10.4.1.4, p.~457); thus $A(\psi^*)=A(\psi)^*$ is not automatic. The counterexample \fid{DET-04}, or already $\left(\begin{smallmatrix}1&1\\0&2\end{smallmatrix}\right)$, proves that nonnegative spectral values and power traces do not establish positive selfadjointness in that quotient metric (\fid{MAS-05}).

\subsection{MAS-BOUND: the boundary Sobolev cap needs a strict inequality}
The norm \eqref{eq:masnorm} permits an independent endpoint check. In the coordinate $u=\log|g|_F$, a Mellin-derivative functional with kernel $\chi(g)|g|_F^{b+it}u^m$ has squared dual norm, up to the finite measure of $C_F^1$,
\begin{equation}\label{eq:boundaryintegral}
 \int_\R e^{2bu-2\alpha|u|}|u|^{2m}(1+u^2)^{-s/2}\,du.
\end{equation}
It is finite for all integers $m\ge0$ if $|b|<\alpha$, and, at $b=\pm\alpha$, is finite exactly when
\begin{equation}\label{eq:strictcap}
 2m<s-1.
\end{equation}
\begin{proof}
The dual of weighted $L^2(w\,du)$ is represented by kernels in $L^2(w^{-1}\,du)$. This gives \eqref{eq:boundaryintegral}. On each tail a negative exponential dominates all polynomial factors. At a boundary one tail has zero exponential exponent; its integrand is comparable to $|u|^{2m-s}$ and is integrable exactly when $2m-s<-1$. Near zero there is no singularity for $m\ge0$. For $|b|>\alpha$ one tail grows exponentially and diverges. This proves necessity as well as sufficiency.
\end{proof}
Consequently the number of permitted derivative orders, capped by a zero's multiplicity $N$, is
\[
 \min\{N,\max(0,\lceil(s-1)/2\rceil)\}.
\]
The printed (10.3.2.4) uses $\min\{N,\lfloor(s-1)/2\rfloor+1\}$. At odd integral $s$ it includes the divergent endpoint; for example $s=3,m=1$ gives a logarithmic divergence. The difference is activated only if $N$ reaches the extra claimed order; this calculation does not assert that a particular zero does so. Theorem 10.3.2.1 allows real $s>1$ without excluding odd integers. For the present comparison one can avoid the issue by choosing strip edges away from the relevant zeros, or nonexceptional $s$. This records an actual correction forced by the displayed norm while leaving the main trace-class comparison explicitly conditional.

\paragraph{Grothendieck framework.}
The arithmetic base is $\Spec\mathcal O_F$, but the constructed objects are analytic Hilbert quotients of adelic functions. The divisor analogy in \S1.2.1 of the draft is with a divisor on a Stein complex line; it is not construction of a constructible $\ell$-adic complex on $\Spec\mathcal O_F$. Cote 1 applies to the asserted analytic trace-class operators. SGA 5 supplies the geometric correspondence analogy, not a theorem identifying this quotient with an arithmetic motive. No crystal, divided-power envelope, EGA or FGA construction of such an identification is furnished by these statements.

\subsection{MAS-10: an ordinary Xi determinant on the actual weighted quotient}
The trace-ideal question for the reciprocal-square multiplier has an affirmative answer. The extra ingredient is a sector estimate in the actual weighted Hilbert norm. It avoids any assumption that the zero modes form a Riesz basis or have uniformly bounded biorthogonal vectors.

Specialize to $F=\mathbf Q$, the trivial character of $C_{\mathbf Q}^1$, center $a=1/2$, width $\alpha=1/2$, and $s>1$. Normalize the compact Haar factor to mass one and write $u=\log|g|$. Thus
\[
 H=L^2(\R,w(u)\,du),\qquad
 w(u)=e^{|u|}(1+u^2)^{s/2},\qquad E=H/V,
\]
where $V$ is the closed centered theta-image. This is the actual Hilbert quotient of \cite[(10.2.2.1)--(10.2.2.2), p.~445]{masullo}, in the indicated character sector.

\textbf{Exact conditional input.} Use \cite[Theorem 10.2.2.1, (10.2.2.4)--(10.2.2.6), p.~446]{masullo} only for its full spectral-annihilator description. Let $\rho$ range over distinct nontrivial zeta zeros, with multiplicity $N_\rho$, and put $z_\rho=\rho-1/2$. Classically
\[
 0<\Re\rho<1,\qquad z_\rho\ne0,
\]
by the Euler product, the functional equation, nonvanishing on $\Re s=1$, and $\xi(1/2)\ne0$; see \cite[\S2.12; \S3.2, pp.~46--47]{titch}. These facts require no RH. Every zero is therefore interior to the chosen strip, so there is no boundary truncation. The cited draft claim says that
\[
 \ell_{\rho,m}(f)=\int_{\R}f(u)e^{z_\rho u}u^m\,du,
 \qquad 0\le m<N_\rho,
\]
annihilate $V$ and have norm-dense linear span in the Hilbert dual $E'$. This Hilbert-dual density is also explained in \cite[Remark 10.2.1.2, pp.~440--441]{masullo}. No nuclearity assertion from Theorem 10.3.2.1 is used. The conclusions below are conditional on the spectral-annihilator claim alone, together with the stated classical zeta facts; \emph{RH is not an input}.

Put $B_0=1/2+s/4$ and choose an auxiliary real parameter $\mu>B_0$, distinct from the center $a=1/2$. On the ambient space define
\[
 g_\mu(v)=\frac{e^{-\mu|v|}}{2\mu},\qquad
 K_\mu f=\int_{\R}g_\mu(v)f(\,\cdot-v)\,dv.
\]
The integral is a strong, vector-valued integral for each $f$, as justified below. Let $K$ be its induced operator on $E$.

\begin{proposition}
The operator $K$ is trace class and $I-\mu^2K$ has a bounded inverse. The operator
\[
 R=K(I-\mu^2K)^{-1}
\]
is trace class, is independent of $\mu$, and has the full jet multiplier $-1/z^2$ at $z=z_\rho$. Reflection $Jf(u)=f(-u)$ descends to a unitary involution of $E$ and commutes with $R$. For $E_+=\ker(J-I)$ and $R_+=R|_{E_+}$,
\begin{equation}\label{eq:masactualXi}
 \det_{E_+}(I-wR_+)=\frac{\Xi(\sqrt w)}{\Xi(0)}.
\end{equation}
Each pair $\{\rho,1-\rho\}$ contributes its full multiplicity $N_\rho$, exactly once. Moreover
\begin{equation}\label{eq:masGreennorm}
 \|K\|_1\le\frac{\mu}{\sqrt{\mu^2-B_0^2}}
 \sum_\rho N_\rho\Re\frac1{\mu^2-z_\rho^2},
\end{equation}
where the sum includes all distinct nontrivial zeros, and every real part in it is positive.
\end{proposition}
\begin{proof}
\emph{Translations and the quotient.}
The function $\log\sqrt w$ is Lipschitz, with almost-everywhere derivative
\[
 b(u)=\tfrac12\operatorname{sgn}(u)+\frac{s u}{2(1+u^2)},
 \qquad |b(u)|\le B_0.
\]
Consequently translations $T_vf(u)=f(u-v)$ satisfy
$\|T_v\|\le e^{B_0|v|}$. They are strongly continuous: first check this on compactly supported smooth functions, and then use their density and the locally uniform bound. The functionals $\ell_{\rho,m}$ are continuous because
$\int_{\R}e^{2\Re(z_\rho)u}|u|^{2m}w(u)^{-1}\,du<\infty$. Substitution gives
\begin{equation}\label{eq:masjettranslation}
 \ell_{\rho,m}(T_vf)=e^{z_\rho v}
 \sum_{j=0}^m\binom{m}{j}v^{m-j}\ell_{\rho,j}(f).
\end{equation}
The complete annihilator description therefore implies $T_vV\subset V$: every jet vanishes on $T_vf$ when $f\in V$, and $V$ is the annihilator of its annihilator. Since
$g_\mu\in L^1(\R,e^{B_0|v|}\,dv)$, the strong integral defining $K_\mu$ is bounded and preserves $V$. Under the isometric identification $E=V^\perp$, its quotient operator is the compression
\[
 K=P_{V^\perp}K_\mu|_{V^\perp}.
\]
No operator-norm Bochner integrability of $v\mapsto T_v$ is asserted or required.

\emph{The weak domain and the sector estimate.}
Both $g_\mu$ and $g_\mu'$ belong to the same exponentially weighted $L^1$ space, and
$g_\mu''=\mu^2g_\mu-\delta_0$ in distributions. Hence $y=K_\mu f$ satisfies
\[
 y,y',y''\in H,\qquad y''=\mu^2y-f.
\]
These statements follow first for compactly supported smooth $f$ by convolution, then for arbitrary $f\in H$ by the weighted convolution bounds and convergence in distributions. Let $U:H\longrightarrow L^2(\R)$ be multiplication by $\sqrt w$ and put $h=Uy$. Then
\[
 h\in H^1(\R),\quad D h:=h'-bh=Uy'\in H^1(\R),
 \quad D^2h=Uy''.
\]
Indeed $h'=\sqrt w\,y'+bh$, and the same calculation applies to $y'$. The weak derivative of $b$ is never used, so the cusp of the weight at zero causes no omitted boundary term. Integration by parts in $H^1$ gives
\begin{align*}
 \langle(\mu^2-D^2)h,h\rangle
 &=\|h'\|_2^2+\mu^2\|h\|_2^2-\|bh\|_2^2\notag\\
 &\quad+\langle h',bh\rangle-\langle bh,h'\rangle.
\end{align*}
For example, $D^*=-\partial_u-b$ on $H^1$, and the identity follows from
$-\langle Dh,D^*h\rangle$; this also justifies it on the stated weak domain. With
$\delta=\mu^2-B_0^2>0$ and $c=B_0/\sqrt\delta$, it follows that
\[
 \Re\langle(\mu^2-D^2)h,h\rangle
 \ge\|h'\|_2^2+\delta\|h\|_2^2,
\]
\[
 |\Im\langle(\mu^2-D^2)h,h\rangle|
 \le2B_0\|h'\|_2\|h\|_2
 \le c\bigl(\|h'\|_2^2+\delta\|h\|_2^2\bigr).
\]
The last inequality is $2XY\le X^2+Y^2$. Since
$(\mu^2-d^2/du^2)y=f$, the number $\langle K_\mu f,f\rangle$ is the conjugate of
$\langle(\mu^2-d^2/du^2)y,y\rangle$. Thus
\begin{equation}\label{eq:masGreensector}
 \Re\langle K_\mu f,f\rangle\ge0,
 \qquad |\Im\langle K_\mu f,f\rangle|
 \le c\Re\langle K_\mu f,f\rangle.
\end{equation}
Compression preserves this numerical-range inequality because
$\langle PK_\mu x,x\rangle=\langle K_\mu x,x\rangle$ for $x\in V^\perp$.
Therefore both $K$ and $K^*$ satisfy the same symmetric sector bound.

\emph{A trace-ideal lemma.}
If a bounded operator $T$ satisfies the sector bound in \eqref{eq:masGreensector} and an orthonormal basis $(e_n)$ satisfies
$\sum_n\Re\langle Te_n,e_n\rangle<\infty$, then $T$ is trace class. To prove this, set
\[
 A=(T+T^*)/2\ge0,\qquad Q=(T-T^*)/(2i).
\]
The diagonal sum says
$\sum_n\|A^{1/2}e_n\|^2<\infty$, so $A^{1/2}$ is Hilbert--Schmidt and $A$ is trace class. The sector bound says $-cA\le Q\le cA$ in quadratic forms. It yields
\[
 Q=A^{1/2}CA^{1/2}
\]
for a bounded selfadjoint $C$ with $\|C\|\le c$. Here is the factorization argument. Define on $\operatorname{Ran}A^{1/2}$ the form
$\beta(A^{1/2}x,A^{1/2}y)=\langle Qx,y\rangle$.
It is well defined: if $x\in\ker A$, the positive operators $cA\pm Q$ have zero quadratic value at $x$ and therefore annihilate $x$, implying $Qx=0$. Cauchy--Schwarz for these two positive forms makes $\beta$ bounded. Its diagonal bound
$|\beta(u,u)|\le c\|u\|^2$ extends to the closure and represents a bounded selfadjoint operator $C$ of norm at most $c$. Extend it by zero on the orthogonal complement. Consequently
\[
 T=A^{1/2}(I+iC)A^{1/2},\qquad
 \|T\|_1\le\|I+iC\|\|A^{1/2}\|_2^2
 \le\sqrt{1+c^2}\,\Tr A.
\]
The trace-class estimate follows either from the Hilbert--Schmidt product inequality or directly by expanding the product in an orthonormal basis and applying Cauchy--Schwarz to the sum of the rank-one norms. This proves the lemma without a prior spectral trace identity.

\emph{Full jets and an orthonormal triangular basis.}
For $|\Re z|<\mu$,
\[
 F_\mu(z):=\int_{\R}g_\mu(v)e^{zv}\,dv
 =\frac1{\mu^2-z^2}.
\]
Integrating \eqref{eq:masjettranslation}, with absolutely convergent exponential-polynomial moments, gives
\[
 K'\ell_{\rho,m}
 =\sum_{j=0}^m\binom{m}{j}
 F_\mu^{(m-j)}(z_\rho)\ell_{\rho,j},
 \qquad \lambda_\rho:=F_\mu(z_\rho)=\frac1{\mu^2-z_\rho^2}.
\]
These eigenvalues need not be real. If $z_\rho=b+it$, set $A_\rho=\mu^2-b^2+t^2>0$. Then
\[
 \Re\lambda_\rho=\frac{A_\rho}{A_\rho^2+4b^2t^2}>0,
 \qquad |\lambda_\rho|\le\frac1{\mu^2-1/4+t^2}.
\]
The classical zero count used in \fid{B-01} consequently gives
$\sum_\rho N_\rho|\lambda_\rho|<\infty$.

Enumerate the distinct zeros and list each finite jet block in increasing jet order. Their Riesz representatives form a total linearly independent sequence in $E$, and every finite prefix spans a $K^*$-invariant space. Applying Gram--Schmidt preserves these finite-prefix spaces and gives an orthonormal basis $(e_n)$ such that, in the paper's conjugate-linear-first convention,
\[
 K^*e_n-\overline{\lambda_n}e_n\in\operatorname{span}(e_1,\ldots,e_{n-1}),
 \qquad \langle e_n,K^*e_n\rangle=\overline{\lambda_n}.
\]
The conjugation is required by the anti-linear Riesz identification. Each $\lambda_\rho$ occurs $N_\rho$ times. No bounds on the original basis vectors or their biorthogonal partners enter this argument. Hence
\[
 \sum_n\Re\langle K^*e_n,e_n\rangle
 =\sum_\rho N_\rho\Re\lambda_\rho<\infty.
\]
The trace-ideal lemma and the sector bound prove $K\in\Sone(E)$ and
\eqref{eq:masGreennorm}. This establishes trace class before using compactness or a trace-equals-eigenvalue theorem; summable eigenvalues of a general nonnormal operator have not been substituted for singular-value control.

\emph{Exact spectrum, inverse and the reciprocal-square multiplier.}
Since $K^*$ is now compact, if $P_n$ projects onto the first $n$ basis vectors, then
$\|(I-P_n)K^*\|\to0$. To see this, approximate $K^*$ in norm by a finite-rank operator and use uniform convergence of $P_n$ on its finite-dimensional range. The invariant flags give a block matrix
\[
 K^*=\begin{pmatrix}T_n&B_n\\0&D_n\end{pmatrix},
 \qquad \|D_n\|\longrightarrow0,
\]
with $T_n$ finite triangular of diagonal $\overline{\lambda_1},\ldots,\overline{\lambda_n}$.
If $q\ne0$ differs from all these diagonal values, choose $n$ with $\|D_n\|<|q|$. Both $T_n-qI$ and $D_n-qI$ are invertible, the latter by a Neumann series, and the explicit upper-triangular block inverse makes $K^*-qI$ invertible. Conversely each diagonal value is an eigenvalue in a finite invariant prefix. For a fixed nonzero eigenvalue, take $n$ past all its occurrences with $\|D_n\|$ less than its modulus. Every generalized eigenvector then lies in that prefix, since the tail block minus the eigenvalue is invertible. Its algebraic multiplicity is exactly the finite triangular multiplicity. There are only finitely many occurrences: equality $\lambda_\rho=\lambda_{\rho'}$ is equivalent to $z_\rho^2=z_{\rho'}^2$, hence to $\rho'=\rho$ or $1-\rho$. Thus the nonzero spectrum of $K$, with algebraic multiplicities, is precisely the $\lambda_\rho$; no additional nonzero spectral part remains.

Since $z_\rho\ne0$, no $\lambda_\rho$ equals $1/\mu^2$, so $I-\mu^2K$ is boundedly invertible. The ideal property now proves
$R=K(I-\mu^2K)^{-1}\in\Sone(E)$. On every finite jet block, rational composition gives the multiplier
\[
 \frac{F_\mu(z)}{1-\mu^2F_\mu(z)}=-\frac1{z^2}.
\]
The identity holds with all jets at $z=z_\rho$, since the denominator is nonzero there; equivalently it is an identity in the corresponding truncated Taylor algebra. Its diagonal is $-z_\rho^{-2}$ with full multiplicity. The same compact triangular argument gives exactly these nonzero spectral values and their algebraic multiplicities. Two choices of $\mu$ act identically on every dual jet, so their difference is annihilated by a dense subspace of $E'$ and is zero. This proves independence of the auxiliary parameter.

\emph{Parity and the determinant.}
Reflection $Jf(u)=f(-u)$ is a unitary involution on $H$, since $w$ is even, and
\[
 \ell_{\rho,m}(Jf)=(-1)^m\ell_{1-\rho,m}(f).
\]
The functional equation gives $N_{1-\rho}=N_\rho$ and $z_{1-\rho}=-z_\rho$. The complete annihilator description therefore yields $JV=V$, so $J$ descends to a unitary involution of $E$. The even Green kernel commutes with reflection, as do $K$ and $R$. Hence $R_+=R|_{E_+}$ is trace class on the orthogonal even summand.

For each pair $\{\rho,1-\rho\}$, reflection exchanges two independent jet spaces of dimension $N_\rho$ with the indicated signs. Its $+1$ eigenspace has dimension $N_\rho$. Applying $(I+J')/2$ to the dense full jet span proves density in the even dual. Thus $R_+$ has precisely one eigenvalue $-z_\rho^{-2}$ per pair, with algebraic multiplicity $N_\rho$. If $\gamma_\rho=-iz_\rho$, then $\gamma_\rho$ is a possibly complex zero of $\Xi$ and $\gamma_\rho^{-2}=-z_\rho^{-2}$. The ordinary trace-class determinant theorem of cote~1 and \fid{B-01} give
\[
 \det_{E_+}(I-wR_+)
 =\prod_{\{\rho,1-\rho\}}
 \left(1+\frac{w}{z_\rho^2}\right)^{N_\rho}
 =\frac{\Xi(\sqrt w)}{\Xi(0)}.
\]
The product is absolutely locally uniformly convergent. This proves every assertion without RH.
\end{proof}

The construction in \fid{MAS-10} is an ordinary Fredholm determinant on the stated arithmetic theta quotient, conditional on its full spectral-annihilator description. It does not assert that $R_+$ is selfadjoint or positive in that quotient metric, and equality of its determinant with the determinant of $B$ does not supply a unitary intertwiner. If, under RH, the heat construction of \fid{MAS-08} is made on this same quotient, its reciprocal-square operator equals $R$ because its full action on the separating dual jets is the same. No unresolved small-time trace-norm condition remains for the ordinary determinant in \eqref{eq:masactualXi}.

\subsection{MAS-07--09: smooth-test obstruction, heat realization and the actual metric}
Retain the fixed setup of \fid{MAS-10}: center $1/2$, strip half-width $\alpha=1/2$, trivial compact character, compact Haar mass one, and
\[
 H=L^2(\R,w(u)\,du),\qquad
 w(u)=e^{|u|}(1+u^2)^{s/2},\quad s>1,\qquad E=H/V.
\]
The full interior jets are indexed by $z_\rho=\rho-1/2$ and their orders below the zero multiplicity. The source input remains the conditional full spectral-annihilator claim of \cite[Theorem 10.2.2.1, (10.2.2.4)--(10.2.2.6), p.~446]{masullo}. The preceding Green construction has already proved $R_+\in\mathcal S_1(E_+)$ and $\det(I-wR_+)=G(w)$ from that claim, without RH or the draft's separate nuclearity theorem. The heat construction below additionally uses the latter theorem and the stipulated RH. The final no-similarity conclusion will again require only the spectral-annihilator input.

\paragraph{MAS-07: no single compact smooth test has the required multiplier.}
No compact smooth test $\phi$ has centered Mellin values $-1/z_\rho^2$ at every nontrivial zero. More generally, the same is true of $1/(c^2-z_\rho^2)$ for any fixed $c\ge0$, including the multiplier of $C$ when $c=1/2$.
\begin{proof}
Compact-character averaging and the fixed centering send $\phi$ to some $h\in C_c^\infty(\R)$, with centered Mellin transform
$F(z)=\int h(u)e^{zu}\,du$. Compact Haar mass and compact radial support justify differentiation under this average. Three integrations by parts, with vanishing boundary terms, give
\[
 |z|^3|F(z)|\le
 \sup_{|\sigma|\le1/2}\int_\R |h'''(u)|e^{\sigma u}\,du<\infty
 \qquad(z\ne0,\ |\Re z|\le1/2).
\]
Classically $|\Re z_\rho|<1/2$, and there are zeros with $|z_\rho|\to\infty$. But
$|z_\rho|^3/|c^2-z_\rho^2|\sim|z_\rho|\to\infty$, a contradiction. The denominators do not vanish at a nontrivial zero, since such a zero is nonreal. This proof needs no RH. Under RH it is exactly the obstruction to the prescribed values $1/(\gamma^2+c^2)$ at $z_\rho=i\gamma$.
\end{proof}
The assertion concerns the designated zero-by-zero multiplier, not an unordered eigenvalue multiset after an unspecified permutation. A different kernel or subsequent functional calculus is not excluded.

\paragraph{MAS-08: trace-class heat operators and the same inverse square.}
For this paragraph invoke the stipulated RH and write $\rho=1/2+i\gamma$, $\gamma\in\R\setminus\{0\}$, with multiplicity $N_\gamma$. Every zero is strictly interior to the fixed strip, so all jets
$\ell_{\gamma,m}(f)=\int f(u)e^{i\gamma u}u^m\,du$, $0\le m<N_\gamma$, are retained. In addition to their finite independence and norm-dense span in $E'$, assume the draft's continuous map
\[
 \mathcal A:C_c^\infty(\R)\longrightarrow\mathcal S_1(E)
\]
obtained from the centered radial sector of \cite[Theorem 10.3.2.1(2), (10.3.2.2), pp.~453--454]{masullo}. Its continuity is in trace norm, not just operator norm.

Let $T_vf(u)=f(u-v)$. The proof in \fid{MAS-10} gives a strongly continuous quotient group with
$\|T_v\|\le e^{B_0|v|}$, $B_0=1/2+s/4$. For any $h$ satisfying
$\int|h(v)|e^{B_0|v|}\,dv<\infty$, define
\[
 \mathcal A(h)f=\int_\R h(v)T_vf\,dv,\qquad
 \|\mathcal A(h)\|\le\int_\R|h(v)|e^{B_0|v|}\,dv.
\]
This is a Hilbert-space Bochner integral for each fixed vector $f$, hence a strong operator integral. No Bochner measurability of $v\mapsto T_v$ in $\mathcal B(E)$ is asserted.
Put
\[
 k_t(v)=(4\pi t)^{-1/2}e^{-v^2/(4t)},\qquad H_t=\mathcal A(k_t)\quad(t>0).
\]
Then $H_t\in\mathcal S_1(E)$, with trace-norm continuity on compact subintervals of $t>0$, and
\[
 H_tH_r=H_{t+r},\qquad H_t\longrightarrow I\ \hbox{strongly as }t\downarrow0,
 \qquad \sup_{0<t\le1}\|H_t\|<\infty.
\]
\begin{proof}
Choose $\chi\in C_c^\infty((-1,1))$ with
$\sum_{n\in\Z}\chi(v-n)=1$. One construction divides a nonnegative bump positive on $[-1/2,1/2]$ by the positive periodic sum of its translates. Put
$q_{n,t}(v)=\chi(v)k_t(v+n)$. Then
$k_t(v)=\sum_n q_{n,t}(v-n)$, pointwise and in $L^1(e^{B_0|v|}\,dv)$.

Continuity of the linear map $\mathcal A$ on the fixed-support Fr\'echet space with support in $[-1,1]$ supplies $C,M$ such that
\[
 \|\mathcal A(q)\|_1\le C\sum_{j=0}^{M}\|q^{(j)}\|_\infty.
\]
For $t$ in any fixed interval $[a,b]\subset(0,\infty)$, Gaussian differentiation gives
\[
 \sum_{j=0}^{M}\|q_{n,t}^{(j)}\|_\infty
 \le C_{a,b}(1+|n|)^M
       \exp\!\left(-\frac{((|n|-1)_+)^2}{4b}\right).
\]
Each derivative is a polynomial times the Gaussian; on $|v|\le1$ this is the displayed bound. Since
$\mathcal A(q(\,\cdot-n))=T_n\mathcal A(q)$, multiplication by $e^{B_0|n|}$ still leaves a summable bound. Thus
$\sum_nT_n\mathcal A(q_{n,t})$ converges in trace norm, uniformly on $[a,b]$. The weighted $L^1$ estimate identifies its sum with the strong integral $H_t$. Each summand is trace-norm continuous in $t$, by the conditional test-map continuity, so their sum is too.

For each fixed $f$, vector-valued Fubini and the group law give
\[
 H_tH_rf=\iint k_t(v)k_r(y)T_{v+y}f\,dy\,dv=H_{t+r}f.
\]
The integrable majorant is
$k_t(v)k_r(y)e^{B_0(|v|+|y|)}\|f\|$; Gaussian convolution gives the last equality. This is vectorwise Fubini, not an operator-valued Bochner assertion about the translation group.
Finally substitute $v=\sqrt t\,y$ in the norm bound. For $0<t\le1$ the integrand is dominated by a fixed Gaussian times $e^{B_0|y|}$. This proves uniform boundedness. Applying the same substitution to
$H_tf-f$ and using strong continuity of $T_v$ proves strong convergence by dominated convergence.
\end{proof}

The complete spectral action, including jets, is
\begin{equation}\label{eq:MasHeatJets}
 H_t'\ell_{\gamma,m}
 =\sum_{j=0}^m\binom mj
 F_t^{(m-j)}(i\gamma)\ell_{\gamma,j},
 \qquad F_t(z)=e^{tz^2}.
\end{equation}
Consequently the nonzero spectrum of $H_t$ consists precisely of $e^{-t\gamma^2}$, with algebraic multiplicity $2N_\gamma$ for each $\gamma>0$.
\begin{proof}
Apply a jet to the strong integral and expand $(u+v)^m$ in the translation formula. Gaussian decay dominates every polynomial moment and the exponential factor. The resulting moments are derivatives of the Gaussian Mellin transform $F_t$, proving \eqref{eq:MasHeatJets}. These finite jet actions obey the product rule $M_FM_G=M_{FG}$ by Leibniz's identity.

If a nonzero eigenvalue $\lambda$ of the compact $H_t$ differs from all displayed values, evaluation on an eigenvector by jets of order zero and then successively higher orders makes every evaluation zero. Their totality makes that vector zero, a contradiction. Conversely an order-zero eigenfunctional at a displayed value makes $H_t-\lambda I$ noninvertible, so that value is in the spectrum.

For a displayed isolated value $\lambda$, take its finite-rank Riesz projection. Its transpose acts on a finite jet block as identity when the block's diagonal value is $\lambda$ and as zero otherwise: restrict the resolvent to that invariant block and integrate its finite triangular matrix. Exactly the two blocks at $\gamma$ and $-\gamma$ match, with total dimension $2N_\gamma$. The image of the dense jet span is their finite-dimensional span, which is closed; continuity makes the entire projection range equal that span. This gives the full algebraic multiplicity. These are generalized eigenspaces, or Riesz subspaces; no diagonalizability is presumed at a multiple zero.
\end{proof}

The even kernels commute with the unitary involution $Jf(u)=f(-u)$ from \fid{MAS-10}. Hence $H_{t,+}=H_t|_{E_+}$ is trace class. On each paired Riesz subspace the dual even part has the independent basis
\[
 \ell_{\gamma,m}+(-1)^m\ell_{-\gamma,m},
 \qquad 0\le m<N_\gamma.
\]
Indeed $J'\ell_{\gamma,m}=(-1)^m\ell_{-\gamma,m}$ interchanges the two blocks, and applying $(I+J')/2$ to the dense jet span proves totality on the even part. Thus the algebraic multiplicity of $e^{-t\gamma^2}$ on $E_+$ is $N_\gamma$.

The heat construction gives the exact identity with the Green operator:
\begin{equation}\label{eq:MasHeatInverse}
 Q=\int_0^1H_t\,dt,\qquad
 R=Q(I-H_1)^{-1},\qquad
 R_+=Q|_{E_+}(I-H_{1,+})^{-1}.
\end{equation}
\begin{proof}
Here $Q$ really is a $\mathcal B(E)$-valued Bochner integral: $t\mapsto H_t$ is norm continuous for $t>0$, has separable norm range there, and is bounded near zero. Assigning any value at $0$ does not change the integral. Its truncations $\int_\varepsilon^1H_tdt$ are trace class, and the omitted interval has operator norm at most a constant times $\varepsilon$, so $Q$ is compact. Since $\gamma\ne0$, the just-proved spectrum excludes $1$ from $\operatorname{Spec}H_1$; hence $I-H_1$ has a bounded inverse.

On every finite jet block the operator $Q(I-H_1)^{-1}$ has analytic multiplier
\[
 \frac{\int_0^1e^{tz^2}\,dt}{1-e^{z^2}}=-\frac1{z^2}
 \quad\hbox{near }z=i\gamma.
\]
The denominator is nonzero there. Integration and this finite jet algebra identity include every derivative up to the full zero multiplicity. This is exactly the multiplier of the bounded Green operator $R$ in \fid{MAS-10}. Their transposes therefore agree on all jets, whose span is dense in $E'$, so the bounded operators agree. Restricting their commuting even parts proves the last identity.
\end{proof}
The trace-class property and ordinary determinant of this $R_+$ were already established by \fid{MAS-10}; the identity also gives $Q=R(I-H_1)\in\mathcal S_1(E)$. No estimate of $\int_0^1\|H_{t,+}\|_1dt$ is needed for these conclusions. For comparison, the scalar trace integral is finite:
\[
 \int_0^1\Tr H_{t,+}\,dt
 =\sum_{\gamma>0}N_\gamma\frac{1-e^{-\gamma^2}}{\gamma^2}<\infty.
\]
At each positive time the Hilbert trace theorem and the proved multiplicities give the trace; Tonelli and the zero count give the integral. Scalar trace integrability alone would not justify a trace-norm integral for these nonpositive operators.

\paragraph{MAS-09: the actual quotient metric rules out selfadjointness.}
Under stipulated RH define
\[
 K(x)=\int_\R\frac{e^{ixu}}{w(u)}\,du.
\]
This is a continuous real even function, tends to zero as $|x|\to\infty$, and satisfies $K(x)>0$ for every real $x$.
\begin{proof}
The reciprocal weight is integrable, giving continuity and decay by the $L^1$ Fourier theorem. Put $\beta=s/2>0$. The Gamma integral gives
\[
 (1+u^2)^{-\beta}
 =\frac1{\Gamma(\beta)}\int_0^\infty r^{\beta-1}e^{-r}e^{-ru^2}\,dr.
\]
For $r>0$, use the Cauchy Fourier identity
$e^{-|u|}=\pi^{-1}\int_\R e^{iyu}(1+y^2)^{-1}dy$. For $u\ne0$ it follows by a single residue at $i$ or $-i$ in the decaying half-plane; at $u=0$ it is the arctangent integral. The Gaussian Fourier integral now gives
\[
 \int_\R e^{ixu}e^{-|u|}e^{-ru^2}\,du
 =\frac1\pi\sqrt{\frac\pi r}
   \int_\R\frac{e^{-(x+y)^2/(4r)}}{1+y^2}\,dy>0.
\]
Fubini here is justified by the integrable Cauchy density and Gaussian. Integrate this strictly positive expression against the positive Gamma measure. Changing the original $r,u$ order is justified by the absolute integrable majorant
$e^{-|u|}(1+u^2)^{-\beta}$. This proves strict positivity. Reality and evenness also follow directly from the even reciprocal weight.
\end{proof}

Use the paper's conjugate-linear-first inner-product convention, so that the actual Riesz representation is $\ell_{\gamma,0}(f)=\langle v_\gamma,f\rangle$, where
$v_\gamma(u)=e^{-i\gamma u}/w(u)$. Since the functional annihilates $V$, this vector belongs to $V^\perp$ and represents it isometrically on the Hilbert quotient. In $E_+$ the even representative is
$c_\gamma(u)=\cos(\gamma u)/w(u)$, and direct integration gives
\begin{equation}\label{eq:MasActualCorrelation}
 \langle c_\gamma,c_\delta\rangle
 =\tfrac12\bigl(K(\gamma-\delta)+K(\gamma+\delta)\bigr)>0.
\end{equation}
But
\[
 R_+^*c_\gamma=\gamma^{-2}c_\gamma,\qquad
 H_{t,+}^*c_\gamma=e^{-t\gamma^2}c_\gamma.
\]
Two distinct positive ordinates give distinct real eigenvalues and nonorthogonal eigenvectors. A selfadjoint operator has orthogonal eigenvectors at distinct real eigenvalues, so neither $R_+$ nor any $H_{t,+}$ is selfadjoint or positive in this metric. More generally, this same argument applies to any bounded multiplier whose real values differ at two positive ordinates. This is a calculation in the specified quotient norm, not a replacement model.

\paragraph{MAS-09: no bounded similarity can repair $R_+$.}
Conditional only on the full spectral-annihilator input of \fid{MAS-10}, $R_+$ is not boundedly similar to a selfadjoint operator.
\begin{proof}
First, if RH fails, choose an off-line nontrivial zero. Its centered parameter $z=\rho-1/2$ has nonzero real part and nonzero imaginary part, since zeta has no real nontrivial zero. Therefore $z^2$, and hence $-1/z^2$, is nonreal. By \fid{MAS-10} this is an eigenvalue of $R_+$. Similarity preserves spectrum, whereas a selfadjoint operator has real spectrum, which is already a contradiction.

It remains to consider RH. If some zero has $N_\gamma\ge2$, the independent even jets
$e_0=\ell_{\gamma,0}+\ell_{-\gamma,0}$ and
$e_1=\ell_{\gamma,1}-\ell_{-\gamma,1}$ satisfy, with $F(z)=-1/z^2$,
\[
 R_+'e_0=\gamma^{-2}e_0,\qquad
 R_+'e_1=\gamma^{-2}e_1+F'(i\gamma)e_0,\qquad
 F'(i\gamma)=\frac{2}{(i\gamma)^3}\ne0.
\]
The full interior jets are essential here. Through the Riesz identification this gives a nontrivial generalized eigenvector for $R_+^*$ at a real eigenvalue. A selfadjoint operator $A$ has none: if
$(A-\lambda)y=x$ and $(A-\lambda)x=0$ with $\lambda\in\R$, then
$\|x\|^2=\langle x,(A-\lambda)y\rangle=\langle(A-\lambda)x,y\rangle=0$.
Bounded similarity preserves this obstruction.

If all zeros are simple, the positive zero count $N(T)\sim T\log T/(2\pi)$ forces distinct positive pairs $\gamma_n,\delta_n\to\infty$ with $|\gamma_n-\delta_n|\to0$. Otherwise a positive eventual lower gap would give $N(T)=O(T)$; equivalently, apply the pigeonhole principle to the successive gaps below $T$. Discreteness ensures pairs with gaps tending to zero escape each fixed compact interval.
By \eqref{eq:MasActualCorrelation},
\[
 \|c_\gamma\|^2=\tfrac12(K(0)+K(2\gamma))\longrightarrow K(0)/2,
 \qquad
 \langle c_{\gamma_n},c_{\delta_n}\rangle\longrightarrow K(0)/2.
\]
Thus their normalized vectors $u_n,v_n$ satisfy
$\langle u_n,v_n\rangle\to1$ and $\|u_n-v_n\|\to0$, while remaining eigenvectors of $R_+^*$ for distinct eigenvalues.

If $R_+^*=SAS^{-1}$ with $A$ selfadjoint and $S$ boundedly invertible, then
$S^{-1}u_n,S^{-1}v_n$ are orthogonal, each of norm at least $1/\|S\|$. Their difference has norm at least $\sqrt2/\|S\|$, contradicting its upper bound
$\|S^{-1}\|\|u_n-v_n\|\to0$. Finally, bounded similarity of $R_+$ to a selfadjoint operator would give one for $R_+^*$ by taking adjoints. This proves the assertion in all cases, without a simplicity assumption.
\end{proof}

Consequently the actual ordinary Fredholm determinant constructed in \fid{MAS-10} can equal that of the designated positive diagonal $B$ while its quotient operator has no bounded similarity to a selfadjoint operator. The obstruction distinguishes the determinant and full algebraic spectrum from the Hilbert metric. A change to an unbounded or nonequivalent completion would be another construction; it is not supplied by this argument. Every Masullo-dependent conclusion remains conditional on the exact draft inputs stated above, with nuclearity needed for the Gaussian trace-class family but not for \fid{MAS-10} or the no-similarity conclusion.

\section{The six programmes after the stipulated RH theorem}
The following are consequences through existing criteria, not assertions that those criteria supplied Alp\"oge's proof. A structural construction ``does not follow'' means that the cited implication gives no such construction; no formal logical independence claim is intended. General-Hecke extensions below remain conditional on Masullo's stated GRH claim.

\subsection{PRG-WEIL: the quadratic form, with its sign specified}
Use $\widetilde g(s)=\int_0^\infty g(x)x^{s-1}dx$ and $g^\sharp(x)=x^{-1}\overline{g(x^{-1})}$. For multiplicative convolution,
\[
 \widetilde{g*g^\sharp}(s)=\widetilde g(s)\overline{\widetilde g(1-\bar s)}.
\]
The explicit formula in the attached Connes--Consani normalization reads
\[
 \widetilde f(0)-\sum_\rho\widetilde f(\rho)+\widetilde f(1)=\sum_v\mathcal W_v(f).
\]
Therefore, for smooth compactly supported $g$ with $\widetilde g(0)=\widetilde g(1)=0$, the Alp\"oge consequence is
\[
 Q(g)=\sum_\rho|\widetilde g(\rho)|^2\ge0,
 \qquad \sum_v\mathcal W_v(g*g^\sharp)=-Q(g)\le0.
\]
This follows by the same Mellin calculation as \fid{MAS-04}; rapid decay and the zero count justify the sum. The converse is Weil's criterion, in the precise version \cite[Appendix C, Proposition C.1, label \texttt{mainprop}]{cc}. It is positivity of a global quadratic form, not positivity of each local term or of a scaling generator. The analogous character-sector statement follows from stipulated Dirichlet GRH with the corresponding completed local factors.

The source's local archimedean result, \fid{PRG-CC-LOCAL}, is already unconditional: \cite[Introduction, Theorem 1, label \texttt{mainthmintro}]{cc} proves
$W_\infty(g*g^*)\ge\Tr(\vartheta(g)\mathbf S\vartheta(g)^*)\ge0$
for its specified support $[2^{-1/2},2^{1/2}]$ and transform vanishings at $i/2$ and $0$. Here $\mathbf S$ is the Sonin-space projection and the notation is the paper's centered normalization. That local compression is real operator structure, but is not an identity with $B$. A geometric realization of the global form by arithmetic correspondences or intersection theory remains an additional target. The curve comparison has exact home in SGA 5, III B, 1.3.1; no EGA construction of the number-field realization is obtained by the inequality alone.

\subsection{PRG-LI / B-03: a positive sequence and an actual paired trace}
Li's criterion is \cite[\S1, Theorem 1]{li}. Its coefficients are
\[
 \lambda_n=\frac1{(n-1)!}\left.\frac{d^n}{ds^n}
 \bigl(s^{n-1}\log\xi(s)\bigr)\right|_{s=1}
 =\sum_\rho\left[1-\left(1-\frac1\rho\right)^n\right],\qquad n\ge1,
\]
where the latter is the symmetric zero sum. All $\lambda_n$ are positive under the stipulated RH. The following checked special form connects this criterion to the designated $B$ and refines SA23 \S4.1(4).

\begin{proposition}
Unconditionally, put $C=B(1+B/4)^{-1}$ and $P_n(x)=2-2T_n(1-x/2)$, where $T_n$ is the Chebyshev polynomial. Then
\begin{equation}\label{eq:litrace}
 \lambda_n=\Tr P_n(C)=\Tr\bigl[2I-2T_n(I-C/2)\bigr]\quad(n\ge1).
\end{equation}
Under RH, $P_n(C)\ge0$ and $\lambda_n>0$.
\end{proposition}
\begin{proof}
The operator $1+B/4$ is invertible: a diagonal zero would require $\gamma_j^2=-1/4$, corresponding to $\rho=0$ or $1$, neither a zero of $\xi$; the diagonal entries tend to one, so their inverses are uniformly bounded. Thus $C$ is trace class, with entries
$x_j=(\gamma_j^2+1/4)^{-1}=[\rho_j(1-\rho_j)]^{-1}$.
For $a=1-1/\rho$, the reflected zero $1-\rho$ gives $a^{-1}$ and
$a+a^{-1}=2-x$. The Chebyshev recurrence gives
$a^n+a^{-n}=2T_n(1-x/2)$. The two Li terms therefore sum to $P_n(x)$.

Since $P_n(0)=0$, $P_n(C)$ is trace class, and the pair sum converges absolutely because $P_n(x)=O_n(x)$ near zero. Pairing preserves the symmetric height cutoff. This proves \eqref{eq:litrace}; the displayed bracket is one polynomial with canceled constant term, not a subtraction of two separately defined infinite traces. Under RH each $x_j$ lies in $(0,4)$, so $T_n(1-x_j/2)\le1$ and $P_n(C)\ge0$. There are arbitrarily small positive $x_j$, and $P_n'(0)=n^2>0$, giving strict positivity of the trace.
\end{proof}
No trace-class operator can have the unpaired spectrum $\{1/\rho\}$, by reciprocal-zero divergence. That narrow intake statement remains correct. A broader assertion that Li coefficients have no trace-class expressions is refuted by \eqref{eq:litrace}. Their positivity supplies no identification of the finite automorphic determinants with this $B$.

\subsection{PRG-NB: the precise closure assertion is now true}
In $L^2((0,\infty),dx)$ put $\rho_a(x)=\{1/(ax)\}$ and $\chi=1_{(0,1]}$. B\'aez-Duarte's strengthened Nyman--Beurling criterion, \cite[Theorem 1.1]{baez}, gives the Alp\"oge consequence
\[
 \chi\in\overline{\operatorname{span}\{\rho_a:a\in\mathbb N\}}^{L^2}.
\]
This is membership of the indicated function, not density of that span in all of $L^2(0,\infty)$. There is no remaining RH gate in this equivalent statement. Prescribed approximants, sharper rates or an arithmetic geometric explanation require their own results; the cited paper already provides constructive approximation information. Cote 1 supplies no identification of these dilations with its determinant operator, and the closure statement itself is functional analysis rather than an EGA/FGA assertion.

\subsection{PRG-DB-FAMILY: de Branges, with the entire shift family}
For every $\varepsilon>0$ set
\[
 E_\varepsilon(z)=\xi(1/2+\varepsilon-iz)=\Xi(z+i\varepsilon),\qquad
 E_\varepsilon^\#(z)=\overline{E_\varepsilon(\bar z)}=\Xi(z-i\varepsilon).
\]
The Alp\"oge consequence is the Hermite--Biehler property
\begin{equation}\label{eq:HB}
 |E_\varepsilon(z)|>|E_\varepsilon^\#(z)|\quad(\Im z>0)
\end{equation}
for \emph{every} $\varepsilon>0$. A numbered family criterion in this convention is \cite[Proposition 1.2; (1.3), (1.7), (1.11)]{suzuki}.
\begin{proof}[Proof of this shifted-$\Xi$ form]
Under RH, $\Xi(z)=\Xi(0)\prod_{\gamma>0}(1-z^2/\gamma^2)$. If $z=x+iy$, $y>0$, each real zero satisfies
$|z+i\varepsilon-\gamma|^2-|z-i\varepsilon-\gamma|^2=4y\varepsilon>0$, and likewise for $-\gamma$. Each paired factor has strictly larger modulus at the upper argument. If the lower-argument product vanishes, the inequality is immediate since the upper argument cannot be a real zero. Otherwise the ratios of partial paired products converge, and fixing one factor bounds every later ratio below a number strictly greater than one. This proves \eqref{eq:HB}. There are no real zeros of $E_\varepsilon$.

Conversely, if all shifts satisfy \eqref{eq:HB} and $\Xi(z_0)=0$ with $\Im z_0>0$, choose $0<\varepsilon<\Im z_0$. Then $E_\varepsilon(z_0-i\varepsilon)=0$ in the upper half-plane, contradicting \eqref{eq:HB}. Conjugation symmetry excludes lower-half-plane zeros as well.
\end{proof}
Each shift consequently defines a de Branges space and a positive reproducing kernel
\[
 K_\varepsilon(w,z)=
 \frac{E_\varepsilon(z)\overline{E_\varepsilon(w)}-
 E_\varepsilon^\#(z)\overline{E_\varepsilon^\#(w)}}{2\pi i(\bar w-z)}.
\]
Positivity here is positivity of all finite kernel Gram matrices. Abstract canonical-system existence follows from de Branges reconstruction; the separate target is a prescribed arithmetic Hamiltonian or kernel, not bare existence. Suzuki's particular arithmetic construction has an explicit parameter range (Theorem 2.3 and \S5). A single large shift is insufficient for RH.

Older sufficient shift-inner-product positivity proposals must be distinguished from this family criterion: Conrey--Li \cite[\S3.1, (3.1)--(3.4)]{conreyli} refutes the proposed $\xi$ choices. Those hypotheses are false, not surviving open consequences of RH. Neither the correct kernel criterion nor those false conditions construct an arithmetic motive or crystal.

\subsection{PRG-CONNES: the RH-equivalent cutoff formula also becomes true}
Connes's local and semilocal trace formulas are established results; the separate global cutoff asymptotic is
\begin{equation}\label{eq:connes}
 \Tr(Q_\Lambda U(h))=2h(1)\log'\Lambda+
 \sum_v\int_{k_v^\times}'\frac{h(u^{-1})}{|1-u|}\,d^\times u+o(1),
\end{equation}
with his cutoffs, measures and test functions \cite[\S VIII, (16), Theorem 5]{connes}. Theorem 5 is first printed in positive characteristic. Its number-field analogue uses the prolate cutoffs in \S VIII, (31)--(37), and the paragraph immediately after (37), which explicitly supplies that analogue. With this qualification, the $\zeta$ sector of \eqref{eq:connes} is true under stipulated RH, and the Dirichlet sectors under stipulated Dirichlet GRH. General-number-field Hecke sectors remain conditional on Masullo here.

Thus the truth of this RH-equivalent asymptotic is not a remaining RH problem. A new direct geometric proof, an identification with \fid{B-01}, and an ordinary trace of the \emph{uncut} scaling action are separate matters. Its relevant positivity is the Weil quadratic distribution, not positivity of the generator. The Lefschetz analogy is with SGA 5's fixed-point traces; the actual adelic noncommutative space is not thereby an \`etale cohomology theory of $\Spec\Z$.

\subsection{PRG-DEN: the cohomology is additional structure}
Deninger's proposed determinant is zeta-regularized:
\[
 \widehat\zeta_K(s)=\prod_{i=0}^2
 \det_\infty\!\left(\frac{s-\Theta}{2\pi}\mid
 H^i_{\mathrm{dyn}}(\overline{\Spec\mathcal O_K},\R)\right)^{(-1)^{i+1}}.
\]
The exact proposal and regularization are \cite[\S3, (2)--(3)]{deninger}; the distributional trace, compact-support version, duality and $K$-theory comparison are (5)--(9). The proposed positive cup/star pairing gives $\Theta+\Theta^*=I$ on degree one, in the common-domain form sense; see pp.~8--9 and the general star-flow relation (11). These are more than zero location.

\paragraph{PRG-DEN-MODEL (what RH alone constructs).}
Using the full zero multiset under RH, on $\ell^2(Z)$ define
\[
 Ae_\rho=\gamma_\rho e_\rho,\quad
 D(A)=\{v:\sum_\rho\gamma_\rho^2|v_\rho|^2<\infty\},\quad
 \Theta=\tfrac12I+iA.
\]
This real diagonal $A$ is selfadjoint: the diagonal inverses of $A\pm i$ are bounded and map onto its domain, or directly its adjoint has the same square-summability domain. Hence $\Theta+\Theta^*=I$ on $D(A)$ and its eigenvalues are the stipulated zeros. This proof establishes a spectral model, not the proposed cohomology. It gives no functorial site, cup product, cycles, duality, or $K$-theory comparison, and no new regularized-determinant theorem. Those remain distinct construction targets. The analogous degree-one Dirichlet spectral statement follows from Alp\"oge; extending to all arithmetic schemes does not follow from RH or Dirichlet GRH. SGA 5 supplies a finite-field model, not an EGA/FGA existence theorem for this number-field object.

\section{PRG-FF: the finite-field dictionary, checked on a curve}
Let $C/\F_q$ be smooth, proper and geometrically connected of genus $g$, with $\ell\nmid q$. Write $V=H^1_{\mathrm{et}}(C_{\overline{\F}_q},\Q_\ell)$ and choose the Frobenius convention with eigenvalues $1$ on $H^0$ and $q$ on $H^2$. Then
\begin{align}
 \#C(\F_{q^n})&=1+q^n-\Tr(F^n\mid V),\label{eq:curve_trace}\\
 Z(C,T)&=\frac{\det(1-TF\mid V)}{(1-T)(1-qT)},\label{eq:curve_det}\\
 -\log\det(1-TF\mid V)&=\sum_{n\ge1}\frac{\Tr(F^n\mid V)}{n}T^n.\label{eq:curve_log}
\end{align}
Exact source locators are SGA 5, III B, Corollary 1.3.1, printed p.~143, and Deligne, Weil I, (1.5.1), (1.5.3), (1.5.4), pp.~275--276 \cite{sga5,weil1}. Equation \eqref{eq:curve_log} is finite linear algebra, so no infinite-dimensional trace theorem is needed. The exterior-power coefficients match cote 1 and SGA 6, V, \S2 and (7.1.1).

Purity, \cite[Theorem (1.6)]{weil1}, says $|\alpha_j|=q^{1/2}$ under each complex embedding. Thus in a complex diagonal realization of the semisimplified eigenvalue multiset, $U=q^{-1/2}F$ is unitary and
\[
 P_1(T)=\det(1-\sqrt q\,T U).
\]
This numerical scaling is not a canonical half Tate twist in the ordinary $\ell$-adic category. Nor is an arbitrary identification of an $\ell$-adic vector space with a complex Hilbert space being asserted. The correct analogue is fixed modulus of eigenvalues, not positive selfadjointness of Frobenius.

\begin{proposition}
For $E/\F_3:y^2=x^3-x$, the curve is smooth, $P_1(T)=1+3T^2$, and $U$ has eigenvalues $i,-i$. No nonzero scalar multiple of its Frobenius can be positive.
\end{proposition}
\begin{proof}
The projective cubic is $Y^2Z-X^3+XZ^2=0$. Its $X$ derivative is $Z^2$ in characteristic three, nonzero on the affine chart; at the sole point at infinity $[0:1:0]$, its $Z$ derivative is one. It is therefore a smooth genus-one curve. For each $x=0,1,2$, $x^3-x=0$, giving three affine points plus infinity. Hence $\#E(\F_3)=4$, the Frobenius trace is $3+1-4=0$, and its determinant on $H^1$ is $3$, since $\Lambda^2H^1\simeq\Q_\ell(-1)$. Thus $P_1(T)=1+3T^2$. The two eigenvalues are $\pm i\sqrt3$. Any nonzero scalar multiple still has opposite nonzero eigenvalues, which cannot both be nonnegative real. Similarity cannot change this obstruction.
\end{proof}
In particular $\det(1-wU)=1+w^2$ and $\Tr U^2=-2$, despite purity. The cup pairing on $H^1$ is alternating with target $\Q_\ell(-1)$; Frobenius is a $q$-similitude. Even that relation alone does not imply purity: $D=\diag(1,q)$ obeys $D^{\mathsf T}JD=qJ$ for $J=\left(\begin{smallmatrix}0&1\\-1&0\end{smallmatrix}\right)$, but its eigenvalue moduli are $1$ and $q$. Direct multiplication proves the claimed identity. Polarization positivity, weight-one purity and positivity of an operator are different statements.

\begin{center}\small
\begin{tabular}{p{0.23\textwidth}p{0.33\textwidth}p{0.33\textwidth}}
\toprule
Feature & $\Xi$ construction & Curve over $\F_q$\\\midrule
Object & Zero-based normal $B$ on $\ell^2$ & Frobenius on finite $H^1_{\mathrm{et}}$\\
Determinant & $\det(1-wB)=G(w)$ & $P_1(T)=\det(1-TF)$\\
Power traces & $\sum_j\gamma_j^{-2n}$ & $1+q^n-\#C(\F_{q^n})$\\
RH statement & $B\ge0$ for this diagonal realization & $|\alpha_j|=\sqrt q$; normalized spectral model unitary\\
Missing geometric datum & Arithmetic realization independent of zero input & Actual sheaf cohomology already supplies $V$, duality and trace\\\bottomrule
\end{tabular}\end{center}

\section{ARC-01--05: repairing the archive argument mathematically}
The relevant theorem on cote 1 is the Hilbert trace/determinant theorem on scans 79--80, with the coefficient domination on scan 79 and the genus-zero argument on scans 111--112. Some of the trace theorem text is struck out. We do not turn that editorial status into a modern publication-priority assertion. The purpose here is to check the mathematical mechanism and specify repairs to the supplied text. Original page images would be needed to call a repair a certified transcription correction.

\subsection{ARC-01: the complete corrected genus-zero lemma}
\begin{lemma}
Let $f(z)=\sum a_nz^n$ be a nonzero entire function and $h(z)=\sum b_nz^n$ an entire function of genus zero with $b_n\ge0$ and $|a_n|\le b_n$. If the nonzero zeros $z_j$ of $f$, counted with multiplicity, satisfy $\sum_j|z_j|^{-1}<\infty$, then $f$ has genus zero.
\end{lemma}
\begin{proof}
A genus-zero product $Cz^m\prod_j(1-z/u_j)$ with $\sum|u_j|^{-1}<\infty$ has minimal exponential type. Indeed, for any $\varepsilon>0$, choose a tail with reciprocal sum less than $\varepsilon/2$. Its modulus is at most $e^{\varepsilon|z|/2}$, while the remaining polynomial is at most $C_\varepsilon e^{\varepsilon|z|/2}$. Thus
\[
 |f(z)|\le h(|z|)\le C_\varepsilon e^{\varepsilon|z|}.
\]
Hadamard factorization \cite[Theorem 3.16, p.~92]{mcmullen} at order at most one and the assumed reciprocal summability give
$f(z)=Ce^{bz}g(z)$ with $C\ne0$ and $g(z)=z^m\prod_j(1-z/z_j)$ of genus zero. If factorization is initially written with genus-one factors, the absolutely convergent sum $\sum1/z_j$ is absorbed into $b$. This standard factorization step is the background theorem invoked on scan 111, line 13.

If $b\ne0$, replace $z$ by $(\bar b/|b|)z$ so that the new exponent is $|b|z$. Choose $\pi/3<\theta<\pi/2$. On $|\arg z|\le\theta$ one has $\Re z\ge\cos\theta\,|z|$. Applying the preceding bound for $f$ with $\varepsilon=|b|\cos\theta/2$ yields $|g(z)|\le A e^{-a|z|}$ on that sector, where $a>0$. The genus-zero bound for $g$ gives $|g(z)|\le B e^{a|z|/3}$ everywhere.

Let $\omega=e^{2\pi i/3}$ and $G_0(z)=g(z)g(\omega z)g(\omega^2z)$. For every direction one of the three rotated arguments has argument in $[-\pi/3,\pi/3]$. All have equal modulus. Therefore $|G_0(z)|\le AB^2e^{-a|z|/3}$ everywhere. Liouville's theorem \cite[Theorem 1.18, p.~19]{mcmullen} makes this entire function constant, and the decay makes the constant zero. A product of three nonzero entire functions is not identically zero. This contradiction forces $b=0$ and proves the lemma.
\end{proof}
If the initial genus-zero majorant has arbitrary complex coefficients, coefficientwise absolute values can be majorized by
$|C|(1+z)^m\prod_j(1+z/|u_j|)$; this justifies the positive-coefficient reduction. The required majorization is radial, $h(|z|)$, not $|h(z)|$. For instance $f=1-z$, $h=1+z$, $z=-1$ disproves the latter inequality. The result applies to arbitrary trace-class Hilbert operators via singular-value determinant domination, not through a positivity hypothesis on the original operator.

\subsection{ARC-02: the order in scan 110}
For $\alpha>1$, $P_\alpha(z)=\prod_{n\ge1}(1-z/n^\alpha)$ has order $1/\alpha$, not $1$.
\begin{proof}
Reciprocal summability gives genus zero. Let $N=\lfloor r^{1/\alpha}\rfloor$. At $z=-r$ the first $N$ factors are at least two, hence $\log M(r)\ge N\log2$. Conversely
\[
 \log M(r)\le \sum_{n\le N}\log(1+r/n^\alpha)+r\sum_{n>N}n^{-\alpha}.
\]
The first sum is at most $N\log(2r)-\alpha\log(N!)=O_\alpha(N)$, using $\log(N!)\ge N\log N-N$ and $r<(N+1)^\alpha$. The integral estimate for the tail also gives $O_\alpha(N)$. Thus $\log M(r)$ is bounded above and below by positive multiples of $r^{1/\alpha}$, proving the order.
\end{proof}
The source says ``sauf erreur,'' so an authorial tentative error is a real possibility; mathematical certainty does not identify the mistaken glyph.

\subsection{ARC-03--04: scan 113 requires a liminf and a prefactor}
For a positive nonincreasing sequence $a_n\to0$ define
\[
 r=\inf\{s>0:a_n=O(n^{-1/s})\},\qquad
 \sigma=\inf\{s>0:\sum_na_n^s<\infty\}.
\]
Then
\begin{equation}\label{eq:liminf}
 \frac1r=\liminf_{n\to\infty}\frac{\log(1/a_n)}{\log n},\qquad r=\sigma,
\end{equation}
with extended reciprocal conventions.
\begin{proof}
If $a_n\le Cn^{-1/s}$, the liminf in \eqref{eq:liminf} is at least $1/s$. Conversely if $1/s$ is strictly less than that liminf, the required bound holds for all sufficiently large $n$ and extends to the finite prefix by a constant. Taking thresholds proves the first assertion. If $a_n=O(n^{-1/s})$, then $\sum a_n^t$ converges for every $t>s$, so $\sigma\le r$. If $\sum a_n^s<\infty$, monotonicity gives
\[
 \lfloor n/2\rfloor a_n^s\le\sum_{j=\lceil n/2\rceil}^n a_j^s\longrightarrow0.
\]
Thus $a_n=o(n^{-1/s})$, proving $r\le\sigma$ and the little-$o$ assertion of the manuscript.
\end{proof}
A counterexample to the printed limsup is explicit. Let $N_k=2^{2^k}$ for $k\ge0$, $a_1=1$, and $a_n=N_k^{-1}$ for $N_k\le n<N_{k+1}=N_k^2$. The sequence decreases to zero. The logarithmic ratio is one at $N_k$ and tends to $1/2$ at the right ends of the blocks, so its limsup is one and liminf is $1/2$. Also $a_n\le n^{-1/2}$, while $a_nn^{1/s}$ is unbounded along block ends for $s<2$. Thus $r=2$, refuting the asserted limsup formula.

For the manuscript's alternative proof that $\sum_n a_n<\infty$ implies $na_n\to0$, suppose $a_{n_i}\ge1/n_i$ after a harmless rescaling and $\varepsilon_i=(n_i-n_{i-1})/n_i$, then for $i\ge2$ summability gives $\sum\varepsilon_i<\infty$, and the correct recurrence is
\[
 n_i=\frac{n_1}{\prod_{j=2}^i(1-\varepsilon_j)}
 \le\frac{n_1}{\prod_{j=2}^\infty(1-\varepsilon_j)}<\infty.
\]
The infinite product is positive because eventually $-\log(1-\varepsilon_j)\le2\varepsilon_j$. This contradicts strict increase of the integers. The prefactor is necessary unless an initial value one is explicitly fixed; an initial zero would instead make the first factor vanish. To prove little-$o$, the subsequence threshold is an arbitrary positive constant, reduced to one by rescaling.

\subsection{ARC-05: norm and domain readings that the formulas force}
For a compact Hilbert operator with singular values $s_j$, the nonzero singular values of $\Lambda^n u$ are the products $s_{i_1}\cdots s_{i_n}$ over strictly increasing indices. This follows first for finite-rank singular-value decompositions and then by finite-rank approximation. For $u$ trace class in the first identity and $u$ Hilbert--Schmidt in the second, this gives the finite identities
\[
 \|\Lambda^n u\|_1=\sum_{i_1<\cdots<i_n}s_{i_1}\cdots s_{i_n}
 =\alpha_n((u^*u)^{1/2}),\qquad
 \|\Lambda^n u\|_2^2=\alpha_n(u^*u).
\]
This proves the readings $\sqrt{u^*u}$ in scan 78, (16), and subscript $2$ in (17). At $n=1$, $u=I_2$, the latter square-root expression is $\sqrt2$, whereas $\|I_2\|_1=2$, so reading subscript $1$ is concretely false. The mixed bound (15) uses nuclear norms; at $n=1$, $u=I_N$, operator norm would give the false inequality $N\le1$. This is a norm-convention clarification, not proof that an omitted subscript was present in the manuscript.

In scan 81, $f$ is defined on the open convex set $U$. Its supporting affine functions satisfy $L_a(x)\le f(x)$ for $x\in U$ and $L_a(a)=f(a)$. For $x\in U$, choosing $a=x$ proves $f(x)=\sup_aL_a(x)$ with attainment. Without an extension the expression $f(x)$ for $x\notin U$ is undefined. Even the natural extension $F=\sup_aL_a$ need not attain its supremum outside $U$: on $U=(0,\infty)$, take $f(x)=1/x$, whose tangents have $L_a(0)=2/a$. Then $F(0)=+\infty$ and no tangent attains it. The required domain in item 3 is therefore $x\in U$, unless an extension and extra attainment hypotheses are supplied.

The expression $e^{z^\omega}$ on scan 110, line 1, is an entire example with the stated order and genus when $\omega$ is a nonnegative integer. For an integer $m>0$, $M(r)=e^{r^m}$ gives order $m$, and the exponential polynomial has genus $m$; $m=0$ gives a nonzero constant of order and genus zero. A nonintegral power has nontrivial monodromy near zero, which exponentiation cannot cancel identically on a punctured neighborhood; a negative integral power gives an essential singularity at zero. Thus the mathematical example needs that qualification, though it may have been implicit in its context.

\subsection{What the correction table does and does not assert}
The table records the product sign on 111/9, rotation on 111/16, sector angle on 112/3, radial majorant on 112/7 and its margin, third root on 112/17, order on 110/2--3, liminf on 113/8 and product indices on 113/22. It also records the unambiguous Hilbert--Schmidt norm on 78/(17), the square root $\sqrt{u^*u}$ on 78/(16), and the domain qualification on scan 81. Contextual issues such as the multiplicative constant and the absorbed $\beta$ factor in 112/6 are explicitly marked as possible implicit normalizations. We do not force a choice between $\sqrt{uu^*}$ and $\sqrt{u^*u}$ in scan 79: their relevant nonzero singular-value data agree, so this mathematics does not resolve that reading. Nor do we call the source's word ``constant'' wrong merely because exponential decay yields the stronger contradiction of a zero constant.

\section{Where these findings live in Grothendieck's frameworks}
\begin{longtable}{>{\raggedright\arraybackslash}p{0.23\textwidth}>{\raggedright\arraybackslash}p{0.40\textwidth}>{\raggedright\arraybackslash}p{0.27\textwidth}}
\toprule
Findings & Exact textual home / comparison & Scope of the relationship\\\midrule\endhead
B-01--03, DET-01--05, DET-08 & Cote 1, scans 78--80, (15)--(23); scans 111--112. SGA 6, V, 2.1, 2.2(b), 2.2.1, (7.1.1). & Analytic trace-class limits and finite exterior-power identities; no arithmetic realization supplied.\\
DET-06--07, DET-10, STEP7-01--02 & Cote 1/79, comparison with $|u|$; SGA 6, V, 2.2(b). & Finite period/Gram/Capelli determinants; a signed determinant is not a positive metric.\\
DET-09, DET-11, STEP6-01--03 & Cote 1/79--80, trace and resolvent; SGA 5, III B, 1.3.1 as the curve analogue of an explicit formula. & Exact scalar prime-sum and Weil-form traces on the zero model; historical arithmetic matrix identification remains separate.\\
TOY-01--02 & SGA 5, III, 4.7--4.8, for a realized sheaf trace; cote 1/79 for the distinct determinant trace. & Finite-field schemes and character sheaves when the weight has that realization; arbitrary weights need no sheaf.\\
MAS-01--10, MAS-BOUND & Cote 1/79--80 and 111--112; SGA 5, III, 4.7 only as geometric analogy. & The full spectral quotient claim gives an actual ordinary Xi determinant; the original metric is not positive. Draft dependence remains; no \'etale cohomology or crystal is identified.\\
PRG-WEIL, PRG-LI, PRG-NB, PRG-DB & Cote 1/79--80 for applicable operator identities; SGA 5, III B, 1.3.1 supplies the curve comparison. & The criteria themselves belong to their cited analytic works; no invented EGA/FGA locator.\\
PRG-CONNES, PRG-DEN & SGA 5, III, 4.7--4.8 and III B, 1.3.1; cote 1/79--80 separates ordinary from regularized determinants. & Proposed arithmetic analogues require their own constructions and adjoint/duality structure.\\
PRG-FF & SGA 5, III B, 1.3.1; Weil I, (1.5.1), (1.5.3)--(1.5.4), (1.6). & Actual $\ell$-adic cohomology, weights, finite determinants and the correct Frobenius convention.\\
ARC-01--05 & Cote 1, scans 78, 81, 110--113, individual lines in the transcription table. & Analytic repairs established; manuscript readings await the images.\\\bottomrule
\end{longtable}
An exact EGA or FGA citation cannot honestly locate a number-field Hilbert positivity construction that those texts do not make. Likewise this prompt produces no new theorem about crystals or divided powers. The operative common language here is exterior powers, traces, determinants, duality and weights. In schemes over $\Z$, local purity of an actual $\ell$-adic sheaf and global RH for its completed $L$-function are distinct assertions; this audit does not extend the stipulated degree-one results to general Hasse--Weil $L$-functions.

\part{Unchecked identifications and explicit limits of the result}
\section{The remaining gate is an identification theorem, not the trace formula}
\paragraph{U-STEP-14 (historical source matching).}
The note's step 1 names a Stepanov determinant, and step 4 names $W(1_S)\ge0$ via a determinant; neither gives a matrix. The full-source exact-name and independent structural searches do not recover a unique declaration for either phrase. Archive metadata and the public mirror history supply no October chronology. The exact candidate Weil-form calculation in \fid{STEP6-01--03} reconstructs the two-indicator test, its scalar prime-sum cross term and its half-plane consequence; it narrows the remaining dictionary to normalization, linear versus quadratic convention, the sign/amplitude and the actual character/window specialization. The step-1 search read actual modular Wronskians, theta evaluation determinants, finite monomial interpolation matrices, directional Hasse derivatives of polynomial minors, and Frobenius/Moore determinants. The step-4 search distinguished finite Pl\"ucker indicator determinants, alternating Weil-named forms, quotient measures, positive Rankin--Selberg integrands and the finite inertia theorem of \fid{DET-10}. Exact locators and read scope are preserved in the evidence notes. Each inspected candidate has a specified finite algebraic or analytic domain; resemblance does not establish the missing historical correspondence. Step 7 is now supported by actual finite minor, exterior-power and period identities, but the note does not identify every determinant with one declaration. This remains an exact-source correspondence gap. The audit-constructed interlacing selector and the actual rank-one period factorization in \fid{STEP7-01--02} further distinguish the finite mechanisms. A full formal proof is not rechecked or rejected on that account.

\paragraph{U-LIMIT (the stronger possible interpretation).}
The literal Fable proposal and the fixed finite polynomial identification are refuted. It remains unproved that a transformed family of Alp\"oge's determinants is a compatible family of compressions or approximants of \fid{B-01}. No impossibility theorem about every such future transformation is asserted. The current sufficient gate is concrete: exhibit the actual operators on a common Hilbert space; establish their Hermitian positivity and trace-norm convergence; and prove their limiting spectral determinant equals $G$ independently of assuming RH. The checked \fid{DET-05} then finishes the desired implication. Alternatively provide a positive arithmetic kernel and a unitary/spectral identification theorem.

The approaches tried were: the supplied note and map; exact entry-file reading; a literal finite Fourier identification (refuted by eigenvalues); path-to-trace identification (refuted and replaced by the rank-one variation); fixed-degree determinant identification (refuted by infinitely many zeros); actual Gram/Capelli/exterior-period source analysis; and the scalar zero-free endpoint trace. A global textual search for \texttt{Fredholm}, \texttt{riemannXi} and \texttt{xiFunction} returns no occurrences, while trace-class automorphic operators do occur. This is supporting search evidence only, not a proof that an unnamed equivalent construction is impossible.

\paragraph{U-MAS (conditional source claims remain conditional).}
The proof of the complete Masullo draft is not a finding of this run. The actual single-operator Fredholm bridge is now proved in \fid{MAS-10}, conditional on its full spectral-jet quotient claim; there is no remaining trace-class gap in that construction. The draft's stated trace-class/arithmetic formulas and general-Hecke GRH consequence retain their conditional labels, the $I$-quantifier qualification, the distinct Chapter 13 statement, and the boundary correction. We do not fill the empty Theorem 13.3.1.2 with an invented assertion. Determining the truth of those deep draft claims is a separate audit from identifying their statements and comparing their operator types.

\paragraph{U-TRANSCRIPTION (images needed).}
The mathematics specifies the corrections; attribution of each difference to the transcription rather than to Grothendieck's draft awaits the original images. The required next step is the second-pass comparison of the listed cote/scan/line entries, not another analytic proof.

All findings in Part I have complete checked proofs or are explicitly attributed conditional applications of cited results. Part II contains no new claimed theorem with an unacknowledged gap. The status file records the exact limits above rather than declaring a universal identification disproved by a failed text search. No three-hour or hourly movement review is claimed before its scheduled time; elapsed-time checkpoints are handled in the dated ledger.

\begin{thebibliography}{99}\raggedright
\bibitem{archive} A.~Grothendieck, unpublished archive, cote 1, supplied transcriptions: scans 78--80, formulas (15)--(23); scan 81, item 3; scan 110, lines 1--12; scan 111, lines 1--23; scan 112, lines 1--23; scan 113, lines 1--23. No manuscript image is part of the supplied ZIP.
\bibitem{groth56} A.~Grothendieck, \emph{La th\'eorie de Fredholm}, Bull. Soc. Math. France \textbf{84} (1956), 319--384. Chapter II, \S4, Theorem 2 (p.~354) and Theorem 4 (p.~359): invertibility and zero multiplicity. \url{https://www.numdam.org/item/10.24033/bsmf.1476.pdf}.
\bibitem{alpoge} L.~Alp\"oge, \emph{Remark}, supplied four-page note dated in the prompt 20 November 2026. Timeline \S2: steps 1, 3, 4 and footnote 3, p.~3; steps 7, 8, 16 and 20, p.~4. \url{https://alpo.ge/remark.pdf}. The theorem convention is that of the task; no Lean replay is claimed.
\bibitem{lean} Alp\"oge Lean sources: supplied \path{ALPOGE_Main.lean}, declarations \texttt{riemannHypothesis}, \texttt{riemannZeta\_eq\_zero\_re}, \texttt{LFunction\_ne\_zero\_of\_half\_lt\_re}, \texttt{LFunction\_eq\_zero\_re}, \texttt{LFunction\_eq\_zero\_re\_of\_isPrimitive}; propositions in \path{ALPOGE_Statements.lean}. Additional definition files and exact declarations are cited in the text and source manifest. Public mirror snapshot: \url{https://github.com/CocoSgt/riemann-hypothesis-lean/releases/tag/v2026-11-20}.
\bibitem{map} \emph{ALPOGE\_LEAN\_MAP\_20261120.md}, prior textual mapping run, Claude Opus 5.5, effort medium, 20 November 2026; \S3 dependency levels and \S4 step table. Textual mapping only; no build claimed by that source.
\bibitem{sa23} \emph{INTAKE\_RESEARCH\_SA23\_20261120.md}, supplied SA23 run, Claude Opus 5.5, effort max as identified in the task; \S1.H and \S4.1(1)--(4), especially \texttt{geo-H-bounds}, \texttt{geo-H-genus}, \texttt{geo-H-lidskii}, \texttt{geo-H-1956-reconstruction}. The present paper checks the stated $\Xi$ consequences and repairs their overbroad interpretations.
\bibitem{fable} \emph{FABLE\_PROGRAM\_SUGGESTIONS\_20261120.md} and \emph{FABLE\_PROGRAM\_SUGGESTIONS\_20261120\_EXPLAINED.md}, supplied Fable 5.2 run, effort max as identified in the task, suggestion 3 in each file. Toy theorem and proposed Fourier/Fredholm and Landau identifications.
\bibitem{masullo} A.~M.~Masullo, \emph{The Music of the Primes}, supplied 940-page draft, dated 20 November 2026 in the task. Exact statements used: Definition 10.1.1.1; Theorems 10.2.2.1, 10.3.1.1, 10.3.2.1, 12.5.1.1, 13.3.1.1; the empty heading 13.3.1.2; Definitions 10.4.1.1, 10.4.1.3, 10.4.3.7; Remark 10.4.1.4; norm formulas (9.2.2.34), (9.2.2.40)--(9.2.2.41); introductory \S\S1.2, 1.4.5--1.4.8, 1.5. All claimed general-Hecke consequences remain conditional on the draft.
\bibitem{titch} E.~C.~Titchmarsh, \emph{The Theory of the Riemann Zeta-Function}, 2nd ed., revised by D.~R.~Heath-Brown, Oxford, 1986. \S2.1, Theorem 2.12 and (2.12.4)--(2.12.5), pp.~29--30; Theorem 9.4, (9.4.3), p.~214. Original book scan inspected at \url{https://sites.math.rutgers.edu/~zeilberg/EM18/TitchmarshZeta.pdf}.
\bibitem{cc} A.~Connes and C.~Consani, \emph{Weil positivity and trace formula, the archimedean place}, supplied source \path{LIT_Connes_Consani_Weil_positivity_archimedean/weil-compo.tex}. Introduction, Theorem 1 (label \texttt{mainthmintro}); Appendix A for the Mellin normalization; Appendix C, Proposition C.1 (label \texttt{mainprop}), equation \texttt{appendweilnegav}. Section-based numbering must be retained.
\bibitem{li} X.-J.~Li, \emph{The positivity of a sequence of numbers and the Riemann hypothesis}, J. Number Theory \textbf{65} (1997), 325--333, \S1, Theorem 1. \url{https://www.sciencedirect.com/science/article/pii/S0022314X97921375}. Statement checked in a reproduction of the original article; no independent re-proof of the converse criterion is claimed.
\bibitem{baez} L.~B\'aez-Duarte, \emph{A strengthening of the Nyman--Beurling criterion for the Riemann hypothesis, 2}, arXiv:math/0205003v1, Theorem 1.1, p.~2. \url{https://arxiv.org/pdf/math/0205003}.
\bibitem{suzuki} M.~Suzuki, \emph{A canonical system of differential equations arising from the Riemann zeta-function}, arXiv:1204.1827v2, 23 September 2016. Proposition 1.2; (1.3), (1.7), (1.11); Theorem 2.3; \S5. \url{https://arxiv.org/pdf/1204.1827}.
\bibitem{conreyli} J.~B.~Conrey and X.-J.~Li, \emph{A note on some positivity conditions related to zeta and $L$-functions}, IMRN (2000), no.~18, 929--940. Theorems 1--2 and \S3.1, (3.1)--(3.4). \url{https://arxiv.org/pdf/math/9812166}.
\bibitem{connes} A.~Connes, \emph{Trace formula in noncommutative geometry and the zeros of the Riemann zeta function}, Selecta Math. (N.S.) \textbf{5} (1999), 29--106. Author text: \S V, Theorem 3; \S VII, Theorem 4; \S VIII, Theorem 5 and (16), (31)--(37), paragraph following (37); Appendix II, Theorem 6. Numbering here follows \url{https://alainconnes.org/wp-content/uploads/selecta.ps-2.pdf}.
\bibitem{deninger} C.~Deninger, \emph{Arithmetic geometry and analysis on foliated spaces}, 23 May 2005. \S3, (2), (3), (5)--(11), Proposition 3.2 and Theorem 3.3; cup/star calculation pp.~8--9. \url{https://swc-math.github.io/dls/DLSDeninger.pdf}.
\bibitem{sga5} A.~Grothendieck and collaborators, \emph{Cohomologie $\ell$-adique et fonctions $L$}, SGA 5, LNM 589. Expos\'e III, Corollary 4.7, (4.7.1), and Corollary 4.8, printed p.~109; Expos\'e III B, Corollary 1.3.1, printed p.~143. Original scans inspected: \url{https://library.slmath.org/nonmsri/sga/sga/5/SGA5-page-121.png} and \url{https://library.slmath.org/nonmsri/sga/sga/5/SGA5-page-155.png}.
\bibitem{sga6} P.~Berthelot, A.~Grothendieck, L.~Illusie and collaborators, \emph{Th\'eorie des intersections et th\'eor\`eme de Riemann--Roch}, SGA 6, LNM 225. Expos\'e V, Definition 2.1, Example 2.2(b), Lemma 2.2.1, (7.1.1); exact text inspected at \url{https://grothendiecksga.com/read/sga6/fr/V-2.html} and \url{https://grothendiecksga.com/read/sga6/fr/V-7.html}.
\bibitem{weil1} P.~Deligne, \emph{La conjecture de Weil. I}, Publ. Math. IHES \textbf{43} (1974), 273--307. (1.5.1), (1.5.3), (1.5.4), Theorem (1.6), pp.~275--276. \url{https://www.numdam.org/item/PMIHES_1974__43__273_0.pdf}.
\bibitem{ramare} O.~Ramar\'e, \emph{Corrigendum to ``Explicit estimates for the summatory function of $\Lambda(n)/n$ from the one of $\Lambda(n)$''}, corrected Lemma 2.1, p.~1, and its integral expression, p.~2. Corrects Lemma 2.2 of Acta Arith. \textbf{159} (2013), no.~2, 113--122. \url{https://ramare-olivier.github.io/Maths/Corrigendum-PsiTildeSimple-Acta-02.pdf}.
\bibitem{mcmullen} C.~McMullen, \emph{Advanced Complex Analysis}, Harvard Math 213a course notes, version dated 16 November 2026. Theorem 1.18, p.~19 (Liouville); Theorem 3.16, p.~92 (finite-order Hadamard factorization). \url{https://people.math.harvard.edu/~ctm/home/text/class/harvard/101/22/html/www/math213a/home/course/course.pdf}.
\end{thebibliography}
\end{document}
